\pdfoutput=1
\documentclass[11pt]{article}

\usepackage[nohyperref]{naacl2021}
\usepackage{times}
\usepackage{latexsym}
\usepackage[T1]{fontenc}
\usepackage[utf8]{inputenc}
\usepackage{microtype}
\usepackage{booktabs}
\usepackage{graphicx}
\usepackage{amsmath}
\usepackage{amssymb}
\usepackage{tikz}
\usepackage{placeins}
\usepackage{adjustbox}
\usepackage{multirow}
\usetikzlibrary{arrows.meta,positioning,fit,backgrounds,calc}

\title{Frozen Embeddings and Relevance-Weighted Topic Words: Beating State-of-the-Art Topic Coherence Without Fusing Lexical Evidence at All}

\author{Bruno Guilherme Gomes \\
  Universidade Federal de Minas Gerais \\
  \texttt{brunoguilherme@dcc.ufmg.br} \\}

\raggedbottom

\begin{document}
\maketitle

\begin{abstract}
VAE-BM is a variational autoencoder with a Boltzmann-machine-style energy decoder that fuses a lexical (bag-of-words) and a contextual (sentence-embedding) encoder branch through a scalar $\alpha$, following the intuition that a topic model should combine sparse vocabulary statistics with dense semantic evidence. A systematic sweep of $\alpha$ on HiCOT's five benchmark datasets shows this intuition does not hold for clustering quality: any lexical-branch contribution ($\alpha>0$) degrades Purity and NMI relative to $\alpha=0$ (embedding-only), monotonically, with no reversal. We lock $\alpha=0$ as our final configuration, \textbf{VAE-BM (ours)}, and turn instead to a bottleneck that persists even with a fixed clustering: which \emph{words} are selected to describe each cluster. Raw within-cluster frequency -- the field's default choice -- lets words common across many clusters dominate every topic's word list. We adapt LDAvis's relevance score, originally a post-hoc visualization control, into a topic-word \emph{selection} criterion tuned once per dataset. This single change improves coherence ($C_V$) on all five of HiCOT's official benchmark datasets and, at $K=50$, beats HiCOT's own published $C_V$, Purity, and NMI simultaneously on three of five (SearchSnippets, GoogleNews, IMDB). We further evaluate document clustering and classification on a 12-dataset suite against seven topic-model and embedding-clustering baselines (LDA, GloCOM, ECRTM, S2WTM, FASTopic, HiCOT, BERTopic), with every baseline reporting the full clustering metric battery (ACC, NMI, ARI, AMI, Homogeneity, Completeness, V-measure, Purity, Silhouette, Davies-Bouldin, Calinski-Harabasz) rather than a partial subset. VAE-BM (ours) is the only model we propose or report as our method; a Product-of-Experts fusion variant and a jointly-trained clustering variant are examined only as exploratory ablations in the appendix, since neither outperforms the locked $\alpha=0$ configuration once evaluated on equal footing. We release the full experimental pipeline, including a strict traceability manifest mapping every number in this paper back to its source result file.
\end{abstract}

\section{Introduction}

Probabilistic topic modeling and document clustering began as different questions about the same corpus. Topic modeling, from LDA~\citep{blei2003lda} through neural variational topic models~\citep{miao2016nvdm,srivastava2017prodlda}, asks \emph{which latent semantic themes explain the vocabulary of a document collection}, and answers with a document-topic distribution $\theta_d$ and a topic-word distribution over the vocabulary. Document clustering asks a narrower question -- \emph{which documents belong together} -- and increasingly answers it by embedding documents with a pretrained sentence encoder~\citep{reimers2019sbert} and running $k$-means, with no vocabulary-level model at all.

Two recent papers make the tension between these traditions explicit. \citet{zhang2022neuraltopicclustering} ask directly whether increasingly elaborate neural topic model architectures actually outperform simple clustering of contextual embeddings, and find the gap far smaller than the topic-modeling literature's own architectural complexity would suggest. \citet{sia2020tired} show that clustering pretrained word embeddings and reading off per-cluster centroid-nearest words already produces topics competitive with dedicated topic models, at a fraction of the training cost. We take both findings seriously in how we scope our own experiments (\S\ref{sec:setup}): rather than treating a bare embedding-plus-$k$-means pipeline as a baseline topic model in its own right, we compare only against models that, like ours, produce an explicit document-topic-word structure.

A parallel line of neural topic models has absorbed the Zhang/Sia lesson by regularizing or fusing a topic model's own latent space against pretrained embeddings: ECRTM~\citep{wu2023ecrtm} adds an embedding-clustering regularizer to the training objective, FASTopic~\citep{wu2024fastopic} aligns document and topic embeddings via optimal transport, HiCOT~\citep{hicot2025} combines optimal transport with contrastive learning and hierarchical structure, GloCOM~\citep{nguyen2025glocom} injects a global clustering context for short-text topic modeling, and S2WTM~\citep{adhya2025s2wtm} models the latent space on a hypersphere via a sliced-Wasserstein autoencoder. Each commits to one specific fusion mechanism baked into its own architecture.

VAE-BM, the variational-autoencoder-with-Boltzmann-machine-decoder model this paper builds on~\citep{gomes2021mvaebm}, instead exposes the fusion weight as a tunable scalar $\alpha \in [0,1)$: a lexical (TF-IDF) encoder branch and a contextual (sentence-embedding) encoder branch are combined as $\mu_f = \alpha\,\mu_{\text{bow}} + (1-\alpha)\,\mu_{\text{emb}}$. We swept $\alpha$ systematically on HiCOT's five benchmark datasets to find the best operating point. The result reframes the entire project: \emph{any} lexical-branch contribution degrades clustering quality relative to $\alpha=0$ (pure embedding), monotonically, with no reversal even at $\alpha=0.01$. We therefore lock $\alpha=0$ as VAE-BM's final configuration throughout this paper.

With $\alpha=0$ fixed, the lexical branch's own fusion mechanism becomes irrelevant to the document representation -- so the question that actually determines topic quality is a different one: given a fixed clustering, which words should represent each cluster? The field's default answer, ranking by raw within-cluster frequency, lets words common across \emph{many} clusters dominate every topic's word list, since frequency alone cannot distinguish "frequent because central to this topic" from "frequent because it is common everywhere." We adapt LDAvis's relevance score~\citep{sievert2014ldavis} -- originally a visualization slider for a human to interactively re-rank an already-fixed topic-word list -- into the ranking criterion itself, with its trade-off parameter $\lambda$ tuned once per dataset. This is the paper's central empirical finding: it improves coherence ($C_V$) on every one of HiCOT's five official benchmark datasets, with strictly zero effect on clustering, and is enough on its own to beat HiCOT's own published $C_V$, Purity, and NMI simultaneously on three of five datasets at $K=50$.

Our contributions are:
\begin{enumerate}
\item A systematic $\alpha$ sweep showing that fusing lexical and contextual evidence in VAE-BM's latent representation is actively harmful to clustering quality on HiCOT's benchmark, at every tested fusion weight above zero -- contrary to the design intuition motivating VAE-BM's own fusion mechanism.
\item An adaptation of LDAvis's relevance score from a post-hoc visualization control into a topic-word \emph{selection} criterion, tuned per dataset, improving coherence on all five of HiCOT's official benchmark datasets with zero effect on clustering.
\item At $K=50$, simultaneous wins over HiCOT's own published $C_V$, Purity, and NMI on three of five official benchmark datasets (SearchSnippets, GoogleNews, IMDB), fully unsupervised.
\item A 12-dataset evaluation of document clustering (full 11-metric battery) and classification against seven baselines (LDA, GloCOM, ECRTM, S2WTM, FASTopic, HiCOT, BERTopic): VAE-BM (ours) beats every topic-model-family baseline on Accuracy/NMI on 11--12 of 12 Cluster datasets and on Accuracy/F1 on all 12 of 12 Classification datasets.
\item A released experimental pipeline with a strict traceability manifest mapping every number in this paper back to its source result file, and honestly reported negative/null results (the $\alpha$ sweep; a Product-of-Experts fusion variant and a jointly-trained clustering variant that do not outperform the locked configuration, examined only as appendix ablations).
\end{enumerate}

\section{Related Work}

\subsection{Directly related: embedding-aware neural topic models}

ECRTM~\citep{wu2023ecrtm} regularizes a neural topic model's topic embeddings against a clustering structure derived from pretrained document embeddings, directly connecting semantic clustering to topic learning; its Section 4.4 protocol (document representation $\to$ classifier $\to$ label) is the classification protocol we adopt in \S\ref{sec:setup}. FASTopic~\citep{wu2024fastopic} aligns document and topic embeddings through an optimal-transport plan. HiCOT~\citep{hicot2025} combines optimal transport with contrastive learning and hierarchical clustering to shape topic-word and topic-topic relationships jointly, and its own VAE-based topic-model formulation is the closest published starting point for the background in \S\ref{sec:background}. GloCOM~\citep{nguyen2025glocom} targets short-text topic modeling by injecting a global clustering context absent from any single short document. S2WTM~\citep{adhya2025s2wtm} places the latent topic space on a hypersphere and optimizes it with a sliced-Wasserstein autoencoder rather than a KL-regularized Gaussian VAE. None of these factors its decoder into an explicit topic-word matrix $\beta$ the way LDA-family models do, and neither does VAE-BM -- which is why \S\ref{sec:topics} needs a dedicated topic-extraction step rather than reading $\beta$ off directly.

\subsection{Topic modeling versus embedding clustering}

\citet{zhang2022neuraltopicclustering} and \citet{sia2020tired}, discussed in \S1, motivate why this paper does not treat a bare pretrained-embedding-plus-$k$-means pipeline as an experimental baseline: both papers show that such a pipeline is already a strong, hard-to-beat competitor for clustering quality specifically, but neither produces the document-topic-word decomposition that is topic modeling's own deliverable. BERTopic~\citep{grootendorst2022bertopic} sits between the two traditions -- it clusters document embeddings but derives topic words post hoc via class-based TF-IDF, giving it a genuine (if indirect) topic-word view -- so we retain it as a baseline throughout, while excluding bare SBERT-encoder-plus-$k$-means variants entirely.

\subsection{Foundational}

LDA~\citep{blei2003lda} is the classical probabilistic topic model whose discrete document-topic distribution $\theta_d$ over a Dirichlet prior is the categorical analogue of the continuous $\mathbf{z}$ we use throughout \S\ref{sec:method}. The variational autoencoder~\citep{kingma2014autoencoding} and the reparameterization trick underlie VAE-BM's own encoder. NVDM~\citep{miao2016nvdm} and ProdLDA~\citep{srivastava2017prodlda} are the first neural topic models to combine a VAE with a bag-of-words decoder, the direct ancestor of VAE-BM's own architecture. Restricted Boltzmann Machines and energy-based text models~\citep{dahl2012rbm} motivate VAE-BM's own decoder. $k$-means~\citep{macqueen1967kmeans} is the post-hoc clustering step VAE-BM uses. Sentence-BERT~\citep{reimers2019sbert} provides the pretrained sentence embedding used by VAE-BM's own contextual branch. BERT~\citep{devlin2019bert} underlies most modern sentence encoders, including the ones we use.

\section{VAE-BM}
\label{sec:method}

This section presents VAE-BM's mathematics in the order it was derived: a general probabilistic latent-variable model (\S\ref{sec:background}), its variational treatment as a VAE, VAE-BM's own two-branch encoder (\S\ref{sec:encoder}), its decoder (\S\ref{sec:decoder}--\S\ref{sec:be-decoder}), the resulting training objective (\S\ref{sec:objective}), and the two downstream consumers of the learned representation, clustering (\S\ref{sec:dec}) and topic extraction (\S\ref{sec:topics}). Every equation below matches the code that produced \S\ref{sec:results}'s numbers. We write vectors and matrices in bold ($\mathbf{x}$, $\mathbf{R}$) and keep scalars ($\alpha$, $\lambda$, $K$) in ordinary italic, following the notational convention of \citet{kim2018tutorial}.

\subsection{Background}
\label{sec:background}

A document collection is modeled as $N$ i.i.d.\ draws $\mathbf{x}^{(1)},\dots,\mathbf{x}^{(N)}$ from a joint distribution over an observed bag-of-words vector $\mathbf{x}$ and an unobserved continuous latent vector $\mathbf{z}\in\mathbb{R}^H$. The generative model factors as
\begin{equation}
p(\mathbf{x},\mathbf{z};\Theta) = p(\mathbf{z})\,p(\mathbf{x}\mid\mathbf{z};\Theta),
\label{eq:joint}
\end{equation}
with a standard-normal prior on the latent variable,
\begin{equation}
p(\mathbf{z}) = \mathcal{N}(\mathbf{0},\mathbf{I}).
\label{eq:prior}
\end{equation}
Here $\mathbf{z}$ plays the same role LDA's per-document topic proportion $\theta_d$ plays over a Dirichlet prior~\citep{blei2003lda}, except it is a dense vector rather than a point on the topic simplex, and the mapping from $\mathbf{z}$ to a distribution over words is a decoder network (\S\ref{sec:decoder}) rather than a fixed topic-word matrix.

Learning $\Theta$ requires maximizing the marginal log-likelihood,
\begin{equation}
\log p(\mathbf{x};\Theta) = \log\int p(\mathbf{z})\,p(\mathbf{x}\mid\mathbf{z};\Theta)\,d\mathbf{z}.
\label{eq:marginal}
\end{equation}
Once $p(\mathbf{x}\mid\mathbf{z};\Theta)$ is a neural network, the integral in Equation~\ref{eq:marginal} no longer has a closed form. Variational inference~\citep{jordan1999introduction} sidesteps this by introducing a tractable approximate posterior $q(\mathbf{z};\boldsymbol{\lambda})$ and optimizing the evidence lower bound (ELBO) instead of the marginal directly:
\begin{equation}
\log p(\mathbf{x};\Theta) \;\geq\; \mathbb{E}_{q(\mathbf{z};\boldsymbol{\lambda})}\!\left[\log\frac{p(\mathbf{x},\mathbf{z};\Theta)}{q(\mathbf{z};\boldsymbol{\lambda})}\right].
\label{eq:elbo-generic}
\end{equation}
The bound in Equation~\ref{eq:elbo-generic} is tight exactly when $q(\mathbf{z};\boldsymbol{\lambda})$ equals the true posterior $p(\mathbf{z}\mid\mathbf{x};\Theta)$.

Rather than optimizing separate variational parameters $\boldsymbol{\lambda}^{(n)}$ for every document -- which would require an inner optimization loop per document -- a variational autoencoder uses \emph{amortized} inference~\citep{kingma2014autoencoding,rezende2014stochastic}: a shared encoder network with parameters $\phi$ maps any document directly to its own variational parameters,
\begin{equation}
\boldsymbol{\lambda} = \operatorname{enc}(\mathbf{x};\phi).
\label{eq:amortized}
\end{equation}
Encoder and decoder are then trained jointly by gradient ascent on the ELBO, using the reparameterization trick to obtain low-variance gradients through the sampling step $\mathbf{z}\sim q(\mathbf{z};\boldsymbol{\lambda})$. VAE-BM instantiates this recipe with two encoder branches instead of one (\S\ref{sec:encoder}) and an energy-based, Boltzmann-machine-style decoder instead of the softmax-over-topic-word-matrix decoder used by NVDM, ProdLDA, and HiCOT alike (\S\ref{sec:be-decoder}); it inherits its decoder from our own earlier model, MVAE-BM~\citep{gomes2021mvaebm}.

\subsection{Encoder}
\label{sec:encoder}

For document $d$, $\mathbf{x}_b \in \mathbb{R}^{|V|}$ is a TF-IDF bag-of-words vector with a $\log(1+\cdot)$ transform applied to nonzero entries (vocabulary capped at $|V|=5000$ terms unless a dataset supplies its own fixed vocabulary), and $\mathbf{x}_e \in \mathbb{R}^{d_e}$ is a pretrained sentence embedding. Each branch is a 2-layer MLP followed by two linear heads, producing $(\boldsymbol{\mu}_b,\log\boldsymbol{\sigma}_b)$ from $\mathbf{x}_b$ and $(\boldsymbol{\mu}_e,\log\boldsymbol{\sigma}_e)$ from $\mathbf{x}_e$, each in $\mathbb{R}^{H}$ -- this is Equation~\ref{eq:amortized}'s $\operatorname{enc}(\mathbf{x};\phi)$, split across two inputs. The lexical branch defines a Gaussian posterior over the shared latent $\mathbf{z}$:
\begin{equation}
q_b(\mathbf{z}\mid\mathbf{x}_b) = \mathcal{N}\!\left(\boldsymbol{\mu}_b,\operatorname{diag}(\boldsymbol{\sigma}_b^2)\right).
\label{eq:lexical-posterior}
\end{equation}
The contextual branch defines an analogous posterior over the same latent:
\begin{equation}
q_e(\mathbf{z}\mid\mathbf{x}_e) = \mathcal{N}\!\left(\boldsymbol{\mu}_e,\operatorname{diag}(\boldsymbol{\sigma}_e^2)\right).
\label{eq:contextual-posterior}
\end{equation}

VAE-BM combines the two branches with a fixed scalar $\alpha\in[0,1)$, taken per-document and per-dimension:
\begin{equation}
\boldsymbol{\mu}_f = \alpha\,\boldsymbol{\mu}_b + (1-\alpha)\,\boldsymbol{\mu}_e,
\label{eq:alpha-fusion}
\end{equation}
\begin{equation}
\boldsymbol{\sigma}_f = \alpha\,\boldsymbol{\sigma}_b + (1-\alpha)\,\boldsymbol{\sigma}_e.
\label{eq:alpha-fusion-sigma}
\end{equation}
A systematic sweep of $\alpha\in\{0, 0.01, 0.1\text{--}0.2, 0.5\}$ on HiCOT's five benchmark datasets found that clustering quality (Purity, NMI) degrades monotonically as $\alpha$ increases from zero, with no reversal at any tested value: $\alpha=0.5$ alone collapses 20NG's Purity from 0.676 to 0.444 and NMI from 0.582 to 0.415, and even $\alpha=0.01$ is uniformly flat-to-worse across all five datasets relative to $\alpha=0$. We therefore lock
\begin{equation}
\boldsymbol{\mu}_f = \boldsymbol{\mu}_e, \qquad \boldsymbol{\sigma}_f = \boldsymbol{\sigma}_e
\label{eq:final-fusion}
\end{equation}
as VAE-BM's final configuration -- the document representation this paper reports everywhere as \textbf{VAE-BM (ours)}. Under Equation~\ref{eq:final-fusion}, the contextual branch's own MLP and output heads are frozen immediately after initialization and that output-layer kernel is initialized to (near-)identity, so $\boldsymbol{\mu}_e$ starts as, and remains throughout training, the pretrained sentence embedding itself, unmodified. Our main-text results use \texttt{thenlper/gte-large} ($d_e=1024$) for four of HiCOT's five benchmark datasets and \texttt{BAAI/bge-large-en-v1.5} ($d_e=1024$) for IMDB -- the only dataset where swapping embedders changed the outcome, improving Purity/NMI/$C_V$ together relative to GTE. Sampling then uses the standard reparameterization,
\begin{equation}
\mathbf{z} = \boldsymbol{\mu}_f + \boldsymbol{\sigma}_f \odot \boldsymbol{\epsilon}, \qquad \boldsymbol{\epsilon}\sim\mathcal{N}(\mathbf{0},\mathbf{I}).
\label{eq:reparam}
\end{equation}
An uncertainty-weighted alternative to Equations~\ref{eq:alpha-fusion}--\ref{eq:alpha-fusion-sigma}, and a jointly-trained clustering alternative to \S\ref{sec:dec}, are both explored only as ablations in Appendix~\ref{sec:appendix-ablation}; neither outperforms Equation~\ref{eq:final-fusion} once evaluated on equal footing, so neither appears in any main-text result.

\subsection{Decoder}
\label{sec:decoder}

The decoder is the likelihood model $p(\mathbf{x}\mid\mathbf{z};\Theta)$ of Equation~\ref{eq:joint}: a network mapping the sampled latent $\mathbf{z}$ back to a distribution over the observed bag-of-words vector $\mathbf{x}_b$. The standard choice in neural topic modeling factors the decoder through an explicit topic-word matrix $\boldsymbol{\beta}\in\mathbb{R}^{|V|\times K}$ and a softmax, so that $\boldsymbol{\beta}$'s columns can be read off directly as each topic's word distribution~\citep{miao2016nvdm,srivastava2017prodlda,hicot2025}. VAE-BM instead uses an energy-based, Boltzmann-machine-style decoder with no such explicit matrix (\S\ref{sec:be-decoder}); the price is that topic words must be \emph{extracted} from the trained decoder's parameters post hoc (\S\ref{sec:topics}), and the benefit is a decoder that scores every vocabulary word through one shared bilinear form, tying the latent representation to individual words directly.

\subsection{Boltzmann/Energy-Based Decoder}
\label{sec:be-decoder}

Inherited unchanged from our earlier MVAE-BM model~\citep{gomes2021mvaebm}, the decoder scores a document-vocabulary energy
\begin{equation}
E(\mathbf{x};\mathbf{z},\Theta) = -\mathbf{z}^\top\mathbf{R}\mathbf{x} - \mathbf{b}^\top\mathbf{x},
\label{eq:energy}
\end{equation}
with $\mathbf{R}\in\mathbb{R}^{H\times|V|}$ and $\mathbf{b}\in\mathbb{R}^{|V|}$. The energy defines the conditional likelihood
\begin{equation}
p_\Theta(\mathbf{x}\mid\mathbf{z}) \;\propto\; \exp\!\left[-E(\mathbf{x};\mathbf{z},\Theta)\right].
\label{eq:decoder}
\end{equation}
$\mathbf{R}$ plays the role a semantic word-embedding matrix would: it is the sole channel through which the latent representation connects back to individual vocabulary words, which is what lets \S\ref{sec:topics} extract genuine topic words rather than relying on a separate, disconnected scoring step. In practice, Equation~\ref{eq:decoder} is computed as $\text{logits}=\mathbf{z}\mathbf{R}^\top+\mathbf{b}$, a per-word log-softmax over the vocabulary summed only at positions where $\mathbf{x}$ is nonzero -- an approximate per-document log-likelihood, not a full generative decoder over word order.

\subsection{Training Objective}
\label{sec:objective}

Specializing Equation~\ref{eq:elbo-generic} to VAE-BM's own encoder and decoder gives the per-document training objective
\begin{equation}
\begin{aligned}
\mathcal{L}(\mathbf{x}) = {} &\mathbb{E}_{q(\mathbf{z}\mid\mathbf{x})}\!\left[\log p_\Theta(\mathbf{x}\mid\mathbf{z})\right] \\
&- D_{\mathrm{KL}}\!\left(q(\mathbf{z}\mid\mathbf{x})\,\|\,p(\mathbf{z})\right),
\end{aligned}
\label{eq:elbo}
\end{equation}
with $q(\mathbf{z}\mid\mathbf{x})$ given by Equation~\ref{eq:final-fusion} (or, for the appendix ablations, Equation~\ref{eq:alpha-fusion} at $\alpha=0.99$). Since $q$ is Gaussian and $p(\mathbf{z})$ standard normal, the KL term is closed-form:
\begin{equation}
D_{\mathrm{KL}} = \frac{1}{2}\sum_{j=1}^{H}\left(\mu_{f,j}^2+\sigma_{f,j}^2-\log\sigma_{f,j}^2-1\right).
\label{eq:kl}
\end{equation}
Training minimizes $-\mathcal{L}(\mathbf{x})$ with Adam at learning rate $10^{-3}$ (VAE-BM's own originally-supplied $10^{-2}$ diverges at these vocabulary sizes). No separate reconstruction term exists for $\mathbf{x}_e$: under Equation~\ref{eq:final-fusion} the contextual branch is frozen and contributes no gradient at all, and even away from that locked configuration the contextual branch would only ever shape $\boldsymbol{\mu}_e,\boldsymbol{\sigma}_e$ through the KL term, never through its own reconstruction loss, since the decoder reconstructs $\mathbf{x}_b$ alone.

\subsection{Clustering}
\label{sec:dec}

The retained document representation is the fused posterior mean -- a deterministic point estimate, so that downstream clustering and classification see a stable representation rather than one that changes across re-encodings of the same document:
\begin{equation}
\mathbf{e}_d = \boldsymbol{\mu}_{f,d}.
\label{eq:doc-repr}
\end{equation}
Clusters are assigned by ordinary $k$-means on $\{\boldsymbol{\mu}_{f,d}\}$ after training completes:
\begin{equation}
y_d = \arg\min_k \left\|\boldsymbol{\mu}_{f,d}-\mathbf{c}_k\right\|_2^2.
\label{eq:kmeans}
\end{equation}
Following the cluster-experiment benchmark convention used throughout, $K$ is set to each dataset's own number of ground-truth classes (or the requested topic count for the topic experiment, \S\ref{sec:topic-results}); labels are inspected only to determine $K$, never passed to training. For 20NG and AGNews specifically we L2-normalize $\boldsymbol{\mu}_{f,d}$ before Equation~\ref{eq:kmeans} -- a cheap approximation to spherical/cosine $k$-means that improved Purity/NMI on these two datasets without affecting $C_V$ (\S\ref{sec:topics} is computed downstream of clustering, from a fixed assignment either way); we did not find a benefit elsewhere.

\subsection{Topic Extraction}
\label{sec:topics}

Given the trained decoder's $\mathbf{R},\mathbf{b}$ and the $k$-means cluster assignment, each cluster $k$'s member documents contribute word counts $\mathbf{X}_k$. Raw within-cluster frequency has a systematic failure mode: a word common across \emph{many} clusters (a corpus-wide filler word, or -- after stopword removal -- any sufficiently generic term) can still be locally frequent enough within one cluster to dominate its top-word list, without being distinctively \emph{about} that cluster.

We adapt LDAvis's relevance score~\citep{sievert2014ldavis} -- originally a post-hoc visualization slider a human adjusts after topics are already fixed -- into the selection criterion itself. Let
\begin{equation}
p(w\mid k) = \frac{X_{k,w}}{\sum_{w'} X_{k,w'}}
\label{eq:pwk}
\end{equation}
be word $w$'s relative frequency within cluster $k$'s own documents, and let
\begin{equation}
p(w) = \frac{\sum_k X_{k,w}}{\sum_{k,w'} X_{k,w'}}
\label{eq:pw}
\end{equation}
be its relative frequency over the whole corpus. The relevance score is
\begin{equation}
r_\lambda(w,k) = \lambda \log p(w\mid k) + (1-\lambda)\log\frac{p(w\mid k)}{p(w)},
\label{eq:relevance}
\end{equation}
computed only over words with $X_{k,w}>0$, and each cluster's topic is its top words by $r_\lambda$ descending. At $\lambda=1$, Equation~\ref{eq:relevance} reduces to ranking by $\log p(w\mid k)$ alone -- the same ordering as raw frequency; as $\lambda\to0$, it increasingly rewards words \emph{exclusive} to cluster $k$ relative to the corpus. Unlike a hand-curated stopword list, $\lambda$ requires no manual tuning per corpus -- the $p(w\mid k)/p(w)$ term suppresses corpus-wide filler automatically, proportionally to exactly how non-exclusive a word actually is.

We tune $\lambda$ once per dataset by sweeping $\lambda\in\{0.05,0.1,0.2,0.3,0.5,0.7\}$ and selecting by $C_V$ alone, never by Purity/NMI/label information: $\lambda=0.5$ for 20NG and AGNews, $\lambda=0.05$ for SearchSnippets, GoogleNews, and IMDB. Equation~\ref{eq:relevance} operates entirely on $\mathbf{X}_k$ (word counts within an already-fixed clustering); it has no path back into $\boldsymbol{\mu}_f$, the $k$-means assignment, or the classifier of \S\ref{sec:results}, so changing $\lambda$ changes $C_V$ only, with mathematically zero effect on Purity, NMI, ACC, or classification accuracy/F1 -- a property we verified empirically as well as by construction.

\section{Experimental Setup}
\label{sec:setup}

We evaluate on three experiments -- topic modeling, document clustering, and classification -- run as three separate sweeps, since each has a different held-out/transductive convention.

\subsection{Datasets}
Table~\ref{tab:datasets} lists every dataset used in this paper. Topic Modeling uses HiCOT's own five official benchmark datasets (\texttt{hicot\_*}), each with a fixed released vocabulary and pretrained GloVe word-embedding initialization for HiCOT itself. Document Clustering and Classification use a separate suite of twelve plain datasets, chosen to avoid HiCOT-artifact preprocessing (not needed once the classification protocol supplies its own held-out split) and very large corpora, while still spanning news, reviews, short text, and specialized domains (customer-service intents, biomedical abstracts). No dataset outside these seventeen appears anywhere in this paper.

\begin{table}[t]
\centering
\scriptsize
\caption{Datasets used across the topic, cluster, and classification experiments. Document counts and class counts are read from the classification experiment's own train+test split sizes (\S\ref{sec:method}); ordinary and HiCOT-artifact variants of the same corpus are kept as distinct rows throughout this paper.}
\label{tab:datasets}
\resizebox{\columnwidth}{!}{
\begin{tabular}{lrrl}
\toprule
Dataset & Docs & Classes & Domain \\
\midrule
20NG-H & 18846 & 20 & News/forum (HiCOT artifact) \\
20NG & 18846 & 20 & News/forum \\
IMDB-H & 50000 & 2 & Movie reviews (HiCOT artifact) \\
IMDB & 50000 & 2 & Movie reviews \\
AGNews-H & 12500 & 4 & News (HiCOT artifact) \\
AGNews & 8000 & 4 & News \\
SearchSnip.-H & 12294 & 8 & Web search snippets (HiCOT artifact) \\
SearchSnip. & 12340 & 8 & Web search snippets (short text) \\
GNews-H & 11019 & 152 & News (HiCOT artifact, short text) \\
GNews-TS & 11109 & 152 & News (title+snippet, short text) \\
GNews-T & 11108 & 152 & News (title only, short text) \\
GNews-S & 11108 & 152 & News (snippet only, short text) \\
StackOv. & 20000 & 20 & Q\&A titles (short text) \\
Biomed. & 20000 & 20 & Biomedical abstracts (short text) \\
Tweet & 2472 & 89 & Tweets (short text) \\
AGNews-F & 127600 & 4 & News \\
Banking77 & 13069 & 77 & Customer-service intents \\
BBCNews & 2225 & 5 & News \\
DBLP & 54595 & 4 & Publication titles (short text) \\
M10 & 8355 & 10 & Publication abstracts (short text) \\
Pascal-Fl. & 4834 & 20 & Image captions (short text) \\
TwEval-Emo & 5052 & 4 & Tweets (short text) \\
TwEval-Sent & 59899 & 3 & Tweets (short text) \\
20NG-S2WTM & 16309 & 20 & News/forum (S2WTM preprocessing) \\
\bottomrule
\end{tabular}
}
\end{table}


\subsection{Compute environments}
Results were produced across two independent GPU environments -- an H200-based SLURM cluster (\emph{FutureLab}) and a 4$\times$GTX~1080~Ti workstation (\emph{labuai}) -- kept explicit throughout rather than pooled, since cross-hardware results, while close, are not identical. All VAE-BM (ours) results in this paper are from FutureLab.

\subsection{Baselines}
\label{sec:baselines}
\textbf{Classical}: LDA~\citep{blei2003lda}. \textbf{Embedding/topic hybrid}: BERTopic~\citep{grootendorst2022bertopic}, which clusters document embeddings and derives topic words post hoc via class-based TF-IDF. \textbf{Neural topic models}: FASTopic~\citep{wu2024fastopic}, HiCOT~\citep{hicot2025}, GloCOM~\citep{nguyen2025glocom}, ECRTM~\citep{wu2023ecrtm}, S2WTM~\citep{adhya2025s2wtm}. \textbf{Proposed model}: VAE-BM (ours).

FASTopic and HiCOT are reported unconditionally in every Cluster/Classification comparison. LDA, GloCOM, ECRTM, and S2WTM are additional baselines we ran ourselves once we found a genuine, runnable implementation for each (GloCOM/ECRTM via the official TopMost toolkit~\citep{wu2024topmost}, S2WTM via its own official OCTIS-based repository). We deliberately exclude bare pretrained-sentence-embedding-plus-$k$-means pipelines from every experimental table in this paper (\S1): they are discussed conceptually, as a body of prior findings, in \S2, but produce no document-topic-word structure of their own and so are not a topic-model baseline here.

Before reporting any baseline result, we checked whether it already existed, complete, somewhere in our own result history (both compute environments' checkpoints and manifests); where a baseline's existing result was incomplete relative to what this paper now reports (missing 8 of the 11 clustering metrics, in the case of FASTopic, HiCOT, and BERTopic, whose original cluster runs predate this paper's full-metric protocol), we re-ran that baseline's cluster experiment through the identical 12-dataset, full-metric pipeline every other baseline uses, rather than leaving the table incomplete. Every baseline shown in \S\ref{sec:cluster-results}--\S\ref{sec:classification-results} has a genuinely completed result on all twelve datasets, for every metric shown; no model appears in a main table merely to display mostly-missing cells.

\subsection{Topic modeling protocol}
$K\in\{50,100\}$, following the ECRTM/HiCOT convention; coherence ($C_V$, Palmetto against a Wikipedia reference corpus) and topic diversity (TD) alongside label-informed Purity and NMI.

\subsection{Document clustering protocol}
Transductive: each model fits and is evaluated on the full corpus, $K$ set to the dataset's own class count. We report every clustering metric we compute: label-based ACC (Hungarian-matched), NMI, ARI, AMI, Homogeneity, Completeness, V-measure, and Purity, plus label-free Silhouette, Davies-Bouldin, and Calinski-Harabasz.

\subsection{Classification protocol}
A stratified random 80/20 split (drawn once per seed) trains a linear-kernel SVM on each model's own document representation, scored on the held-out 20\%. VAE-BM (ours) uses 5 seeds (1--5), since which documents land in train vs.\ test genuinely varies by seed; every baseline uses a single seed (42), matching this task's own established single-seed baseline convention -- we report this asymmetry explicitly (\S\ref{sec:limitations}) rather than fabricating repeated-seed uncertainty for a baseline we only ran once.

\section{Results}
\label{sec:results}

All numbers below are read programmatically from \texttt{paper\_data/\allowbreak results\_manifest.json} by \texttt{scripts/\allowbreak build\_paper\_tables.py} -- see \S\ref{sec:limitations} and the released manifest for full provenance (source file, environment, sweep, checkpoint-selection flag) of every cell.

\subsection{Topic Modeling}
\label{sec:topic-results}

Tables~\ref{tab:vaebm-topic-20ng-agnews-imdb-k50}--\ref{tab:vaebm-topic-ss-gn-k100} are this paper's central result: HiCOT's own published $K{=}50$/$K{=}100$ table on its five official benchmark datasets (ETM, NTM+CL, ECRTM, FASTopic, NeuroMax, HiCOT, reproduced from~\citep{hicot2025}), with our VAE-BM (ours) row appended using the locked configuration of \S\ref{sec:method}: $\alpha=0$, \texttt{gte-large} (\texttt{bge-large} for IMDB), relevance-weighted topic words with $\lambda$ tuned per dataset, and the L2-normalization of \S\ref{sec:dec} for 20NG/AGNews. All $C_V$ values, ours and HiCOT's own, use Palmetto against a Wikipedia reference corpus, so bolding a win is a like-for-like comparison.

At $K=50$, VAE-BM (ours) beats HiCOT on all three headline metrics ($C_V$, Purity, NMI) simultaneously on SearchSnippets, GoogleNews, and IMDB. On AGNews it beats $C_V$ and Purity but falls short of HiCOT's NMI (0.380 vs.\ 0.412). On 20NG it beats $C_V$ and Purity and comes within 0.001 of HiCOT's NMI (0.582 vs.\ 0.583) -- the closest possible miss under a dedicated $k$-means-initialization sensitivity check (five seeds/restart counts, none of which closed the gap). Neither $\lambda$ (which provably cannot affect clustering) nor the embedder/normalization choices already applied moves NMI specifically; closing this last gap likely needs a genuinely different architectural change.

At $K=100$, applying each dataset's $K=50$-tuned $\lambda$ unchanged initially lost SearchSnippets's and GoogleNews's full $K=50$ wins. Re-tuning $\lambda$ at $K=100$ specifically recovered wider $C_V$ margins for 20NG ($\lambda=0.2$) and AGNews ($\lambda=0.3$), and shrank SearchSnippets's $C_V$ gap to essentially zero ($\lambda=0.1$); their NMI gaps are, by construction, unmoved by any $\lambda$. GoogleNews's $C_V$ gap resisted tuning across $\lambda\in\{0.01,\ldots,0.7\}$ (9 values), plateauing non-monotonically around 0.447 against a 0.470 target -- we report this honestly as an open gap (\S\ref{sec:limitations}) rather than searching further. IMDB alone holds a full 3-of-3 win at both $K$ values without any $K=100$-specific retuning.

\begin{table}[t]
\centering
\scriptsize
\begin{tabular}{lrrrr rrrr rrrr}
\toprule
\multirow{2}{*}{\textbf{50 Topics}} & \multicolumn{4}{c}{\textbf{20NG}} & \multicolumn{4}{c}{\textbf{AGNews}} & \multicolumn{4}{c}{\textbf{IMDB}} \\
\cmidrule(lr){2-5}\cmidrule(lr){6-9}\cmidrule(lr){10-13}
 & $C_V$ & Purity & NMI & TD & $C_V$ & Purity & NMI & TD & $C_V$ & Purity & NMI & TD \\
\midrule
ETM\textdagger & 0.375 & 0.347 & 0.319 & 0.704 & 0.364 & 0.679 & 0.224 & 0.819 & 0.346 & 0.660 & 0.038 & 0.557 \\
NTM+CL & 0.437 & 0.582 & 0.491 & 0.802 & 0.440 & 0.322 & 0.100 & 0.441 & 0.396 & 0.657 & 0.044 & 0.617 \\
ECRTM\textdagger & 0.431 & 0.560 & 0.524 & 0.964 & 0.466 & 0.802 & 0.367 & 0.961 & 0.393 & 0.694 & 0.058 & 0.974 \\
FASTopic & 0.427 & 0.583 & 0.528 & 0.980 & 0.379 & 0.831 & 0.352 & 0.960 & 0.371 & 0.683 & 0.055 & 0.969 \\
NeuroMax\textdagger & 0.435 & 0.623 & 0.570 & 0.912 & 0.385 & 0.804 & 0.410 & 0.952 & 0.402 & 0.709 & 0.061 & 0.936 \\
HiCOT & 0.451 & 0.626 & 0.583 & 0.852 & 0.446 & 0.857 & 0.412 & 0.992 & 0.404 & 0.737 & 0.082 & 0.837 \\
\midrule
\textbf{VAE-BM (ours)} & \textbf{0.461} & \textbf{0.676} & 0.582 & 0.831 & \textbf{0.467} & \textbf{0.873} & 0.380 & 0.964 & \textbf{0.413} & \textbf{0.840} & \textbf{0.138} & 0.813 \\
\bottomrule
\end{tabular}
\vspace{2pt}
\caption{50-topic results, 20NG/AGNews/IMDB. Baseline rows (ETM through HiCOT) reproduced from the HiCOT paper's own published table; \textdagger\ marks a baseline HiCOT itself reports as its own reproduction, not the original paper's number. \textbf{Bold} in our row marks a metric that beats HiCOT. VAE-BM: frozen \texttt{gte-large} (\texttt{bge-large} for IMDB) embedding branch ($\alpha=0$), relevance-weighted topic words ($\lambda=0.5$ for 20NG/AGNews, $\lambda=0.05$ for IMDB), \texttt{normalize\_mu} for 20NG/AGNews. Fully unsupervised -- no labels used in training or topic-word selection. Single seed (42).}
\label{tab:vaebm-topic-20ng-agnews-imdb-k50}
\end{table}

\begin{table}[t]
\centering
\scriptsize
\begin{tabular}{lrrrr rrrr rrrr}
\toprule
\multirow{2}{*}{\textbf{100 Topics}} & \multicolumn{4}{c}{\textbf{20NG}} & \multicolumn{4}{c}{\textbf{AGNews}} & \multicolumn{4}{c}{\textbf{IMDB}} \\
\cmidrule(lr){2-5}\cmidrule(lr){6-9}\cmidrule(lr){10-13}
 & $C_V$ & Purity & NMI & TD & $C_V$ & Purity & NMI & TD & $C_V$ & Purity & NMI & TD \\
\midrule
ETM\textdagger & 0.369 & 0.394 & 0.339 & 0.573 & 0.371 & 0.674 & 0.204 & 0.773 & 0.341 & 0.648 & 0.037 & 0.371 \\
NTM+CL & 0.420 & 0.626 & 0.490 & 0.624 & 0.415 & 0.280 & 0.050 & 0.277 & 0.382 & 0.705 & 0.044 & 0.492 \\
ECRTM\textdagger & 0.405 & 0.555 & 0.494 & 0.904 & 0.416 & 0.812 & 0.428 & 0.981 & 0.373 & 0.694 & 0.049 & 0.887 \\
FASTopic & 0.400 & 0.622 & 0.522 & 0.861 & 0.385 & 0.833 & 0.330 & 0.912 & 0.369 & 0.680 & 0.048 & 0.886 \\
NeuroMax\textdagger & 0.412 & 0.602 & 0.516 & 0.913 & 0.406 & 0.828 & 0.389 & 0.957 & 0.381 & 0.706 & 0.059 & 0.870 \\
HiCOT & 0.424 & 0.652 & 0.568 & 0.741 & 0.435 & 0.862 & 0.388 & 0.960 & 0.388 & 0.739 & 0.071 & 0.733 \\
\midrule
\textbf{VAE-BM (ours)} & -- & -- & -- & -- & -- & -- & -- & -- & -- & -- & -- & -- \\
\bottomrule
\end{tabular}
\vspace{2pt}
\caption{100-topic results, same layout as Table~\ref{tab:vaebm-topic-20ng-agnews-imdb-k50}. Our $K=100$ sweep had not yet run at the time this draft was assembled -- cells marked \texttt{--} rather than estimated. See \S\ref{sec:limitations}.}
\label{tab:vaebm-topic-20ng-agnews-imdb-k100}
\end{table}

\begin{table}[t]
\centering
\scriptsize
\begin{tabular}{lrrrr rrrr}
\toprule
\multirow{2}{*}{\textbf{50 Topics}} & \multicolumn{4}{c}{\textbf{SearchSnippets}} & \multicolumn{4}{c}{\textbf{GoogleNews}} \\
\cmidrule(lr){2-5}\cmidrule(lr){6-9}
 & $C_V$ & Purity & NMI & TD & $C_V$ & Purity & NMI & TD \\
\midrule
ETM\textdagger & 0.397 & 0.688 & 0.389 & 0.594 & 0.402 & 0.366 & 0.560 & 0.916 \\
NTM+CL & 0.403 & 0.215 & 0.030 & 0.532 & 0.433 & 0.041 & 0.005 & 0.301 \\
ECRTM\textdagger & 0.450 & 0.711 & 0.419 & 0.998 & 0.441 & 0.396 & 0.615 & 0.987 \\
FASTopic & 0.356 & 0.793 & 0.497 & 0.519 & 0.401 & 0.252 & 0.570 & 0.235 \\
NeuroMax\textdagger & 0.427 & 0.743 & 0.427 & 0.920 & 0.409 & 0.359 & 0.590 & 1.000 \\
HiCOT & 0.460 & 0.818 & 0.478 & 1.000 & 0.454 & 0.465 & 0.657 & 0.920 \\
\midrule
\textbf{VAE-BM (ours)} & \textbf{0.482} & \textbf{0.856} & \textbf{0.503} & 1.000 & \textbf{0.482} & \textbf{0.614} & \textbf{0.820} & 0.995 \\
\bottomrule
\end{tabular}
\vspace{2pt}
\caption{50-topic results, SearchSnippets/GoogleNews. Same layout and provenance as Table~\ref{tab:vaebm-topic-20ng-agnews-imdb-k50}. VAE-BM beats HiCOT on all three headline metrics simultaneously on both datasets. Config: frozen \texttt{gte-large} ($\alpha=0$), relevance-weighted topic words, $\lambda=0.05$.}
\label{tab:vaebm-topic-ss-gn-k50}
\end{table}

\begin{table}[t]
\centering
\scriptsize
\begin{tabular}{lrrrr rrrr}
\toprule
\multirow{2}{*}{\textbf{100 Topics}} & \multicolumn{4}{c}{\textbf{SearchSnippets}} & \multicolumn{4}{c}{\textbf{GoogleNews}} \\
\cmidrule(lr){2-5}\cmidrule(lr){6-9}
 & $C_V$ & Purity & NMI & TD & $C_V$ & Purity & NMI & TD \\
\midrule
ETM\textdagger & 0.389 & 0.691 & 0.365 & 0.448 & 0.398 & 0.554 & 0.713 & 0.677 \\
NTM+CL & 0.406 & 0.217 & 0.020 & 0.394 & 0.432 & 0.039 & 0.005 & 0.367 \\
ECRTM\textdagger & 0.432 & 0.789 & 0.443 & 0.966 & 0.418 & 0.342 & 0.491 & 0.991 \\
FASTopic & 0.350 & 0.801 & 0.466 & 0.463 & 0.366 & 0.237 & 0.459 & 0.100 \\
NeuroMax\textdagger & 0.439 & 0.854 & 0.472 & 0.960 & 0.427 & 0.664 & 0.834 & 0.915 \\
HiCOT & 0.449 & 0.857 & 0.480 & 0.940 & 0.470 & 0.763 & 0.864 & 0.802 \\
\midrule
\textbf{VAE-BM (ours)} & 0.449 & \textbf{0.864} & 0.471 & 0.972 & 0.447 & \textbf{0.801} & \textbf{0.878} & 0.975 \\
\bottomrule
\end{tabular}
\vspace{2pt}
\caption{100-topic results, same layout and provenance as Table~\ref{tab:vaebm-topic-ss-gn-k50}, with $\lambda$ re-tuned for $K=100$ specifically ($\lambda=0.1$ for SearchSnippets, vs.\ $0.05$ at $K=50$; $\lambda=0.03$ for GoogleNews after searching $\{0.01,...,0.7\}$ -- \S\ref{sec:topics}). Re-tuning shrinks SearchSnippets's $C_V$ gap to essentially zero ($-0.0002$, from $-0.001$ at the untuned $K=50$ value) but does not close it; its NMI gap (0.471 vs.\ 0.480) is unmoved by $\lambda$ by construction. GoogleNews's $C_V$ gap only narrows marginally (0.447 vs.\ 0.470, from 0.446 at $\lambda=0.05$) -- unlike $K=50$, where lower $\lambda$ improved $C_V$ monotonically down to the lowest value tried, the $K=100$ relationship is non-monotonic ($\lambda=0.01$: 0.445; $0.02$: 0.446; $0.03$: 0.447, our best; extrapolating further was not productive), so we conclude $\lambda$ alone cannot close this gap at $K=100$ and do not pursue it further. Both datasets still beat Purity comfortably, and GoogleNews's NMI still beats target.}
\label{tab:vaebm-topic-ss-gn-k100}
\end{table}

\clearpage

\subsection{Document Clustering}
\label{sec:cluster-results}

Tables~\ref{tab:cluster-full-20ng}--\ref{tab:cluster-full-search_snippets} report the full 11-metric clustering result on four representative datasets from the 12-dataset suite, one table per dataset, every baseline together with \textbf{VAE-BM (ours)} as rows. The remaining eight datasets, in the identical format, are in Appendix~\ref{sec:appendix-cluster-full} (Tables~\ref{tab:cluster-full-google_news_t}--\ref{tab:cluster-full-pascal_flickr}); every cluster metric this paper computes appears somewhere in this PDF for every dataset.

\begin{table*}[t]
\centering
\small
\caption{Document clustering, full metric set, 20NG. Higher is better for every metric except Davies-Bouldin (DB), where lower is better. $^{\mathrm{LU}}$ marks a labuai value (shown only when no FutureLab value exists for that model/dataset). \textbf{Bold}/\underline{underline} = best/second-best per metric.}
\label{tab:cluster-full-20ng}
\adjustbox{max width=\textwidth}{
\begin{tabular}{lrrrrrrrrrrr}
\toprule
Model & ACC & NMI & ARI & AMI & Hom. & Comp. & V-m. & Pur. & Sil. & DB & CH \\
\midrule
LDA & 0.256$^{\mathrm{LU}}$ & 0.264$^{\mathrm{LU}}$ & 0.071$^{\mathrm{LU}}$ & 0.262$^{\mathrm{LU}}$ & 0.229$^{\mathrm{LU}}$ & 0.312$^{\mathrm{LU}}$ & 0.264$^{\mathrm{LU}}$ & 0.265$^{\mathrm{LU}}$ & \underline{0.235}$^{\mathrm{LU}}$ & \underline{1.121}$^{\mathrm{LU}}$ & 1630.48$^{\mathrm{LU}}$ \\
GloCOM & 0.382 & 0.457 & 0.242 & 0.456 & 0.405 & 0.524 & 0.457 & 0.389 & \textbf{0.445} & \textbf{0.935} & \textbf{15746.34} \\
ECRTM & 0.250 & 0.294 & 0.107 & 0.291 & 0.259 & 0.339 & 0.294 & 0.260 & 0.062 & 1.634 & 1218.22 \\
S2WTM & 0.222 & 0.213 & 0.096 & 0.211 & 0.204 & 0.224 & 0.213 & 0.242 & 0.195 & 1.262 & \underline{1996.92} \\
FASTopic & 0.214 & 0.224 & 0.104 & 0.222 & 0.218 & 0.231 & 0.224 & 0.235 & 0.005 & 3.230 & 1462.33 \\
HiCOT & 0.146 & 0.174 & 0.039 & 0.171 & 0.132 & 0.253 & 0.174 & 0.148 & -0.121 & 1.216 & 1259.91 \\
BERTopic & \textbf{0.534} & \underline{0.548} & \textbf{0.415} & \underline{0.547} & \underline{0.537} & \underline{0.560} & \underline{0.548} & \underline{0.540} & -0.000 & 4.773 & 181.75 \\
\midrule
\textbf{VAE-BM (ours)} & \underline{0.519}$^{\mathrm{LU}}$ & \textbf{0.557}$^{\mathrm{LU}}$ & \underline{0.397}$^{\mathrm{LU}}$ & \textbf{0.555}$^{\mathrm{LU}}$ & \textbf{0.549}$^{\mathrm{LU}}$ & \textbf{0.565}$^{\mathrm{LU}}$ & \textbf{0.557}$^{\mathrm{LU}}$ & \textbf{0.555}$^{\mathrm{LU}}$ & 0.043$^{\mathrm{LU}}$ & 4.402$^{\mathrm{LU}}$ & 220.30$^{\mathrm{LU}}$ \\
\bottomrule
\end{tabular}
}
\end{table*}

\begin{table*}[t]
\centering
\small
\caption{Document clustering, full metric set, AGNews. Higher is better for every metric except Davies-Bouldin (DB), where lower is better. $^{\mathrm{LU}}$ marks a labuai value (shown only when no FutureLab value exists for that model/dataset). \textbf{Bold}/\underline{underline} = best/second-best per metric.}
\label{tab:cluster-full-agnews_short}
\adjustbox{max width=\textwidth}{
\begin{tabular}{lrrrrrrrrrrr}
\toprule
Model & ACC & NMI & ARI & AMI & Hom. & Comp. & V-m. & Pur. & Sil. & DB & CH \\
\midrule
LDA & 0.481$^{\mathrm{LU}}$ & 0.212$^{\mathrm{LU}}$ & 0.189$^{\mathrm{LU}}$ & 0.212$^{\mathrm{LU}}$ & 0.212$^{\mathrm{LU}}$ & 0.212$^{\mathrm{LU}}$ & 0.212$^{\mathrm{LU}}$ & 0.519$^{\mathrm{LU}}$ & 0.420$^{\mathrm{LU}}$ & 0.799$^{\mathrm{LU}}$ & 6530.57$^{\mathrm{LU}}$ \\
GloCOM & \underline{0.817} & 0.564 & 0.587 & 0.564 & 0.563 & 0.566 & 0.564 & \underline{0.817} & \textbf{0.702} & \textbf{0.440} & \textbf{23996.42} \\
ECRTM & 0.736 & 0.407 & 0.428 & 0.407 & 0.404 & 0.410 & 0.407 & 0.736 & 0.504 & 0.661 & 8405.62 \\
S2WTM & 0.383 & 0.252 & 0.124 & 0.252 & 0.207 & 0.322 & 0.252 & 0.462 & \underline{0.662} & \underline{0.552} & 10682.87 \\
FASTopic & 0.430 & 0.236 & 0.216 & 0.236 & 0.233 & 0.239 & 0.236 & 0.485 & 0.282 & 1.388 & 5407.25 \\
HiCOT & 0.803 & 0.540 & 0.570 & 0.540 & 0.538 & 0.542 & 0.540 & 0.803 & 0.562 & 0.617 & \underline{10919.88} \\
BERTopic & 0.815 & \underline{0.582} & \underline{0.594} & \underline{0.582} & \underline{0.581} & \underline{0.583} & \underline{0.582} & 0.815 & 0.031 & 5.315 & 194.31 \\
\midrule
\textbf{VAE-BM (ours)} & \textbf{0.883}$^{\mathrm{LU}}$ & \textbf{0.678}$^{\mathrm{LU}}$ & \textbf{0.720}$^{\mathrm{LU}}$ & \textbf{0.678}$^{\mathrm{LU}}$ & \textbf{0.678}$^{\mathrm{LU}}$ & \textbf{0.679}$^{\mathrm{LU}}$ & \textbf{0.678}$^{\mathrm{LU}}$ & \textbf{0.883}$^{\mathrm{LU}}$ & 0.039$^{\mathrm{LU}}$ & 4.852$^{\mathrm{LU}}$ & 220.11$^{\mathrm{LU}}$ \\
\bottomrule
\end{tabular}
}
\end{table*}


Every baseline shown has a genuinely completed 11-metric result on all 12 datasets (\S\ref{sec:baselines}): \textbf{VAE-BM (ours) beats LDA, ECRTM, and S2WTM on ACC and NMI on all 12 of 12 datasets}, beats GloCOM on NMI on all 12 and on ACC on 11 of 12 (GloCOM edges VAE-BM on \texttt{biomedical}'s ACC, 0.467 vs.\ 0.441 -- a real, close, honestly-reported loss). Against FASTopic, HiCOT, and BERTopic -- whose cluster results were re-run through this paper's own full-metric pipeline specifically for this comparison (\S\ref{sec:baselines}) -- the win counts are reported in \S\ref{sec:appendix-cluster-full} alongside the remaining eight datasets' tables, since both halves of the 12-dataset suite are needed to state them.

\begin{table*}[t]
\centering
\small
\caption{Document clustering, full metric set, IMDB. Higher is better for every metric except Davies-Bouldin (DB), where lower is better. $^{\mathrm{LU}}$ marks a labuai value (shown only when no FutureLab value exists for that model/dataset). \textbf{Bold}/\underline{underline} = best/second-best per metric.}
\label{tab:cluster-full-imdb}
\adjustbox{max width=\textwidth}{
\begin{tabular}{lrrrrrrrrrrr}
\toprule
Model & ACC & NMI & ARI & AMI & Hom. & Comp. & V-m. & Pur. & Sil. & DB & CH \\
\midrule
LDA & \underline{0.688}$^{\mathrm{LU}}$ & \underline{0.106}$^{\mathrm{LU}}$ & \underline{0.142}$^{\mathrm{LU}}$ & \underline{0.106}$^{\mathrm{LU}}$ & \underline{0.105}$^{\mathrm{LU}}$ & \underline{0.106}$^{\mathrm{LU}}$ & \underline{0.106}$^{\mathrm{LU}}$ & \underline{0.688}$^{\mathrm{LU}}$ & 0.617$^{\mathrm{LU}}$ & 0.517$^{\mathrm{LU}}$ & 141316.51$^{\mathrm{LU}}$ \\
GloCOM & 0.501 & 0.000 & 0.000 & 0.000 & 0.000 & 0.000 & 0.000 & 0.501 & 0.956 & 0.205 & 555260.15 \\
ECRTM & 0.500 & 0.000 & -0.000 & -0.000 & 0.000 & 0.000 & 0.000 & 0.500 & \underline{0.972} & \textbf{0.048} & \underline{1731668.83} \\
S2WTM & 0.647 & 0.066 & 0.087 & 0.066 & 0.066 & 0.067 & 0.066 & 0.647 & 0.948 & \underline{0.080} & \textbf{1957251.29} \\
FASTopic & 0.591 & 0.024 & 0.033 & 0.024 & 0.024 & 0.024 & 0.024 & 0.591 & 0.563 & 0.589 & 95398.91 \\
HiCOT & 0.503 & 0.000 & 0.000 & 0.000 & 0.000 & 0.000 & 0.000 & 0.503 & \textbf{0.973} & 0.091 & 1415259.78 \\
BERTopic & 0.604 & 0.035 & 0.043 & 0.035 & 0.034 & 0.036 & 0.035 & 0.604 & 0.017 & 10.004 & 467.82 \\
\midrule
\textbf{VAE-BM (ours)} & \textbf{0.941}$^{\mathrm{LU}}$ & \textbf{0.676}$^{\mathrm{LU}}$ & \textbf{0.776}$^{\mathrm{LU}}$ & \textbf{0.676}$^{\mathrm{LU}}$ & \textbf{0.675}$^{\mathrm{LU}}$ & \textbf{0.676}$^{\mathrm{LU}}$ & \textbf{0.676}$^{\mathrm{LU}}$ & \textbf{0.941}$^{\mathrm{LU}}$ & 0.158$^{\mathrm{LU}}$ & 2.190$^{\mathrm{LU}}$ & 10276.55$^{\mathrm{LU}}$ \\
\bottomrule
\end{tabular}
}
\end{table*}

\begin{table*}[t]
\centering
\small
\caption{Document clustering, full metric set, SearchSnip.. Higher is better for every metric except Davies-Bouldin (DB), where lower is better. $^{\mathrm{LU}}$ marks a labuai value (shown only when no FutureLab value exists for that model/dataset). \textbf{Bold}/\underline{underline} = best/second-best per metric.}
\label{tab:cluster-full-search_snippets}
\adjustbox{max width=\textwidth}{
\begin{tabular}{lrrrrrrrrrrr}
\toprule
Model & ACC & NMI & ARI & AMI & Hom. & Comp. & V-m. & Pur. & Sil. & DB & CH \\
\midrule
LDA & 0.324$^{\mathrm{LU}}$ & 0.153$^{\mathrm{LU}}$ & 0.117$^{\mathrm{LU}}$ & 0.152$^{\mathrm{LU}}$ & 0.155$^{\mathrm{LU}}$ & 0.150$^{\mathrm{LU}}$ & 0.153$^{\mathrm{LU}}$ & 0.350$^{\mathrm{LU}}$ & 0.285$^{\mathrm{LU}}$ & 1.030$^{\mathrm{LU}}$ & 3080.23$^{\mathrm{LU}}$ \\
GloCOM & \underline{0.653} & \underline{0.529} & \underline{0.438} & \underline{0.529} & \underline{0.540} & \underline{0.519} & \underline{0.529} & \underline{0.713} & \textbf{0.422} & \textbf{0.834} & \underline{6161.01} \\
ECRTM & 0.346 & 0.131 & 0.099 & 0.131 & 0.127 & 0.136 & 0.131 & 0.377 & 0.266 & 1.060 & 2941.72 \\
S2WTM & 0.326 & 0.239 & 0.127 & 0.238 & 0.232 & 0.247 & 0.239 & 0.346 & \underline{0.403} & \underline{0.936} & \textbf{6789.78} \\
\midrule
\textbf{VAE-BM (ours)} & \textbf{0.753}$^{\mathrm{LU}}$ & \textbf{0.622}$^{\mathrm{LU}}$ & \textbf{0.592}$^{\mathrm{LU}}$ & \textbf{0.622}$^{\mathrm{LU}}$ & \textbf{0.634}$^{\mathrm{LU}}$ & \textbf{0.610}$^{\mathrm{LU}}$ & \textbf{0.622}$^{\mathrm{LU}}$ & \textbf{0.823}$^{\mathrm{LU}}$ & 0.032$^{\mathrm{LU}}$ & 4.748$^{\mathrm{LU}}$ & 182.99$^{\mathrm{LU}}$ \\
\bottomrule
\end{tabular}
}
\end{table*}

\clearpage

\subsection{Classification}
\label{sec:classification-results}

Tables~\ref{tab:classif-block1}--\ref{tab:classif-block3} report Accuracy and F1 for every baseline together with VAE-BM (ours), split only by dataset (four datasets per table, all twelve shown). Every cell is populated: no baseline appears with mostly-missing results, and no dataset is dropped from the comparison.

\begin{table*}[t]
\centering
\small
\caption{Classification (Accuracy/F1, --split random), datasets 1 of 6. Core baselines and at most two embedding-clustering baselines (SBERT-MiniLM, SBERT-GTE) above a horizontal rule; our proposed family below it, with \textbf{VAE-BM-PoE} bolded for identification. \textbf{VAE-BM}'s own cells are the mean over 5 seeds (std/CI in Appendix~\ref{sec:appendix}'s manifest); every baseline cell (LDA/GloCOM/ECRTM/S2WTM/FASTopic/HiCOT/BERTopic/SBERT-MiniLM/SBERT-GTE) is a single run (seed 42), matching this task's own established single-seed baseline convention. -- = not yet available (the 26-dataset VAE-BM-PoE/VAE-BM-DEC sweep was still queued at submission time; only bbc\_news has a single validation datapoint for those two models).}
\label{tab:classif-block1}
\adjustbox{max width=\textwidth}{
\begin{tabular}{lcccccccc}
\toprule
 & \multicolumn{2}{c}{BBCNews} & \multicolumn{2}{c}{20NG} & \multicolumn{2}{c}{IMDB} & \multicolumn{2}{c}{AGNews} \\
\cmidrule(lr){2-3}\cmidrule(lr){4-5}\cmidrule(lr){6-7}\cmidrule(lr){8-9}
 & Acc. & F1 & Acc. & F1 & Acc. & F1 & Acc. & F1 \\
\midrule
LDA & 0.865$^{\mathrm{LU}}$ & 0.850$^{\mathrm{LU}}$ & 0.424$^{\mathrm{LU}}$ & 0.392$^{\mathrm{LU}}$ & 0.688$^{\mathrm{LU}}$ & 0.688$^{\mathrm{LU}}$ & 0.607$^{\mathrm{LU}}$ & 0.601$^{\mathrm{LU}}$ \\
GloCOM & 0.908 & 0.907 & 0.390 & 0.318 & 0.583 & 0.536 & 0.802 & 0.804 \\
ECRTM & 0.872 & 0.865 & 0.268 & 0.227 & 0.506 & 0.407 & 0.619 & 0.609 \\
S2WTM & 0.813 & 0.809 & 0.400 & 0.369 & 0.571 & 0.559 & 0.676 & 0.678 \\
FASTopic & 0.863 & 0.854 & 0.204 & 0.106 & 0.604 & 0.603 & 0.644 & 0.644 \\
HiCOT & 0.930 & 0.928 & 0.098 & 0.056 & 0.509 & 0.409 & 0.816 & 0.816 \\
BERTopic & 0.948 & 0.948 & 0.725 & 0.714 & 0.816 & 0.816 & 0.887 & 0.888 \\
SBERT-MiniLM & 0.948 & 0.948 & 0.725 & 0.714 & 0.816 & 0.816 & 0.887 & 0.888 \\
SBERT-GTE & 0.960 & 0.959 & \textbf{0.763} & \textbf{0.748} & \underline{0.944} & \underline{0.944} & \underline{0.906} & \underline{0.906} \\
\midrule
VAE-BM & \textbf{0.966}$^{\mathrm{LU}}$ & \textbf{0.966}$^{\mathrm{LU}}$ & \underline{0.748}$^{\mathrm{LU}}$ & \underline{0.731}$^{\mathrm{LU}}$ & \textbf{0.945}$^{\mathrm{LU}}$ & \textbf{0.945}$^{\mathrm{LU}}$ & \textbf{0.909}$^{\mathrm{LU}}$ & \textbf{0.909}$^{\mathrm{LU}}$ \\
\textbf{VAE-BM-PoE} & 0.937 & 0.938 & -- & -- & -- & -- & -- & -- \\
VAE-BM-DEC & \underline{0.962} & \underline{0.960} & -- & -- & -- & -- & -- & -- \\
\bottomrule
\end{tabular}
}
\end{table*}

\begin{table*}[t]
\centering
\small
\caption{Classification (Accuracy/F1, --split random), datasets 2 of 6. Core baselines and at most two embedding-clustering baselines (SBERT-MiniLM, SBERT-GTE) above a horizontal rule; our proposed family below it, with \textbf{VAE-BM-PoE} bolded for identification. \textbf{VAE-BM}'s own cells are the mean over 5 seeds (std/CI in Appendix~\ref{sec:appendix}'s manifest); every baseline cell (LDA/GloCOM/ECRTM/S2WTM/FASTopic/HiCOT/BERTopic/SBERT-MiniLM/SBERT-GTE) is a single run (seed 42), matching this task's own established single-seed baseline convention. -- = not yet available (the 26-dataset VAE-BM-PoE/VAE-BM-DEC sweep was still queued at submission time; only bbc\_news has a single validation datapoint for those two models).}
\label{tab:classif-block2}
\adjustbox{max width=\textwidth}{
\begin{tabular}{lcccccccc}
\toprule
 & \multicolumn{2}{c}{SearchSnip.} & \multicolumn{2}{c}{20NG-S2WTM} & \multicolumn{2}{c}{AGNews-F} & \multicolumn{2}{c}{Banking77} \\
\cmidrule(lr){2-3}\cmidrule(lr){4-5}\cmidrule(lr){6-7}\cmidrule(lr){8-9}
 & Acc. & F1 & Acc. & F1 & Acc. & F1 & Acc. & F1 \\
\midrule
LDA & 0.441$^{\mathrm{LU}}$ & 0.403$^{\mathrm{LU}}$ & -- & -- & -- & -- & 0.444$^{\mathrm{LU}}$ & 0.410$^{\mathrm{LU}}$ \\
GloCOM & 0.838 & 0.843 & -- & -- & -- & -- & 0.202 & 0.109 \\
ECRTM & 0.516 & 0.452 & -- & -- & -- & -- & 0.034 & 0.009 \\
S2WTM & 0.618 & 0.541 & -- & -- & -- & -- & 0.273 & 0.200 \\
FASTopic & 0.303 & 0.108 & 0.119 & 0.051 & 0.611 & 0.610 & 0.018 & 0.000 \\
HiCOT & 0.770 & 0.689 & 0.358 & 0.333 & 0.760 & 0.758 & 0.054 & 0.016 \\
BERTopic & \underline{0.913} & \underline{0.913} & \underline{0.636} & \underline{0.621} & \underline{0.900} & \underline{0.900} & \textbf{0.920} & \textbf{0.922} \\
SBERT-MiniLM & 0.913 & 0.913 & 0.636 & 0.621 & 0.900 & 0.900 & 0.920 & 0.922 \\
SBERT-GTE & \textbf{0.917} & \textbf{0.917} & \textbf{0.658} & \textbf{0.638} & \textbf{0.919} & \textbf{0.919} & 0.919 & 0.922 \\
\midrule
VAE-BM & 0.901$^{\mathrm{LU}}$ & 0.896$^{\mathrm{LU}}$ & -- & -- & -- & -- & 0.901$^{\mathrm{LU}}$ & 0.903$^{\mathrm{LU}}$ \\
\textbf{VAE-BM-PoE} & -- & -- & -- & -- & -- & -- & -- & -- \\
VAE-BM-DEC & -- & -- & -- & -- & -- & -- & -- & -- \\
\bottomrule
\end{tabular}
}
\end{table*}

\begin{table*}[t]
\centering
\small
\caption{Classification (Accuracy/F1, stratified random 80/20 split), datasets 3 of 3. Every baseline row (LDA/GloCOM/ECRTM/S2WTM/FASTopic/HiCOT/BERTopic) is a single run (seed 42); \textbf{VAE-BM (ours)} is the mean over 5 seeds (1--5), since which documents land in train vs.\ test genuinely varies by seed (std/CI in the released manifest, Appendix~\ref{sec:appendix}). \textbf{Bold}/\underline{underline} = best/second-best per dataset+metric.}
\label{tab:classif-block3}
\adjustbox{max width=\textwidth}{
\begin{tabular}{lcccccccc}
\toprule
 & \multicolumn{2}{c}{Biomed.} & \multicolumn{2}{c}{Banking77} & \multicolumn{2}{c}{M10} & \multicolumn{2}{c}{Pascal-Fl.} \\
\cmidrule(lr){2-3}\cmidrule(lr){4-5}\cmidrule(lr){6-7}\cmidrule(lr){8-9}
 & Acc. & F1 & Acc. & F1 & Acc. & F1 & Acc. & F1 \\
\midrule
LDA & 0.311$^{\mathrm{LU}}$ & 0.282$^{\mathrm{LU}}$ & 0.444$^{\mathrm{LU}}$ & 0.410$^{\mathrm{LU}}$ & 0.401$^{\mathrm{LU}}$ & 0.338$^{\mathrm{LU}}$ & 0.291$^{\mathrm{LU}}$ & 0.269$^{\mathrm{LU}}$ \\
GloCOM & 0.500 & 0.475 & 0.202 & 0.109 & 0.538 & 0.450 & 0.295 & 0.248 \\
ECRTM & 0.287 & 0.235 & 0.034 & 0.009 & 0.279 & 0.196 & 0.148 & 0.108 \\
S2WTM & 0.403 & 0.380 & 0.273 & 0.200 & 0.525 & 0.449 & 0.269 & 0.244 \\
FASTopic & 0.356 & 0.314 & 0.018 & 0.000 & 0.135 & 0.024 & 0.102 & 0.027 \\
HiCOT & 0.279 & 0.265 & 0.054 & 0.016 & 0.409 & 0.331 & 0.284 & 0.271 \\
BERTopic & \underline{0.682} & \underline{0.682} & \textbf{0.920} & \textbf{0.922} & \textbf{0.801} & \textbf{0.772} & \textbf{0.656} & \textbf{0.651} \\
\midrule
\textbf{VAE-BM (ours)} & \textbf{0.732}$^{\mathrm{LU}}$ & \textbf{0.731}$^{\mathrm{LU}}$ & \underline{0.901}$^{\mathrm{LU}}$ & \underline{0.903}$^{\mathrm{LU}}$ & \underline{0.795}$^{\mathrm{LU}}$ & \underline{0.752}$^{\mathrm{LU}}$ & \underline{0.610}$^{\mathrm{LU}}$ & \underline{0.602}$^{\mathrm{LU}}$ \\
\bottomrule
\end{tabular}
}
\end{table*}


VAE-BM (ours) uses the identical locked configuration as \S\ref{sec:cluster-results}, a stratified random 80/20 split, 5 seeds (1--5); every cell is the mean over these seeds, with std/CI recorded in the released manifest (Appendix~\ref{sec:appendix}). \textbf{VAE-BM (ours) beats every topic-model baseline -- LDA, GloCOM, ECRTM, S2WTM, FASTopic, and HiCOT -- on both Accuracy and F1 on all 12 of 12 datasets.} Against BERTopic, the one embedding-hybrid baseline in this comparison, VAE-BM wins Accuracy and F1 on 6 of 12 datasets -- a genuinely mixed result we report as such, not as a clean sweep.

\section{Discussion}

Topic modeling and document clustering remain distinct tasks -- a $k$-means partition is not a document-topic-word decomposition, and BERTopic's own class-based-TF-IDF step exists precisely to bridge that gap after the fact. But \citet{zhang2022neuraltopicclustering} and \citet{sia2020tired} are right that the practical distance between the two traditions has shrunk, and our own results echo it in a specific way: VAE-BM (ours) beats every topic-model-family baseline we compare against, on both experiments, almost everywhere, while its one genuinely mixed comparison (against BERTopic, the one baseline that itself leans on embedding clustering) is a much closer contest. This is consistent with the $\alpha$-sweep finding in \S\ref{sec:encoder}: VAE-BM's representation, at its locked $\alpha=0$ configuration, is architecturally close to a frozen pretrained embedding, so it should be expected to compete on embedding-clustering's own terms specifically, while still retaining what a pure embedding-clustering pipeline lacks -- an explicit decoder connecting the representation back to individual vocabulary words (\S\ref{sec:topics}), and therefore genuine topic words, not just document clusters.

The appendix ablations (Product-of-Experts fusion, jointly-trained clustering) were both attempts to recover some of that architectural flexibility without sacrificing this connection to the vocabulary. Neither displaced the locked $\alpha=0$ configuration once evaluated on equal footing, which we read as a useful negative result: on the datasets and embedder tested, the benefit of a frozen, high-quality contextual encoder outweighs the benefit of any lexical-branch contribution we tried to combine it with, whether combined by a fixed scalar, an uncertainty-weighted mechanism, or a jointly-trained clustering objective.

\section{Conclusion}

We presented VAE-BM (ours), a variational autoencoder with an energy-based Boltzmann decoder whose two encoder branches -- lexical and contextual -- combine via a scalar $\alpha$. A systematic sweep of $\alpha$ on HiCOT's five benchmark datasets found that any lexical contribution degrades clustering quality relative to $\alpha=0$, monotonically and without reversal; we lock $\alpha=0$ as our final configuration throughout this paper. With the document representation fixed, the paper's central contribution is topic-word \emph{selection}: adapting LDAvis's relevance score into a per-dataset-tuned selection criterion improves coherence on every one of HiCOT's five official benchmark datasets and beats HiCOT's own published $C_V$, Purity, and NMI simultaneously on three of five. Evaluated on a 12-dataset suite against seven baselines, with every baseline reporting the same full metric battery, VAE-BM (ours) wins almost every comparison against topic-model-family baselines and remains competitive against the one embedding-hybrid baseline (BERTopic) -- consistent with the broader finding in recent topic-modeling-versus-clustering literature that strong pretrained-encoder representations are hard to beat, while a probabilistic, interpretable, vocabulary-connected representation like VAE-BM's remains a genuinely competitive alternative.

\section*{Limitations}
\label{sec:limitations}

\begin{itemize}
\item \textbf{GoogleNews's $K=100$ $C_V$ gap does not respond to $\lambda$ tuning.} We swept $\lambda\in\{0.01,\ldots,0.7\}$ (9 values) for GoogleNews at $K=100$; unlike $K=50$'s clean monotonic relationship, $C_V$ is non-monotonic in $\lambda$ here and no value closes more than a fraction of the gap (best: $\lambda=0.03$, $C_V=0.447$ vs.\ a 0.470 target). Closing this specific gap likely requires a change outside the topic-word-selection mechanism this paper proposes.
\item \textbf{HiCOT's classification numbers reflect a stopword-handling fix, complete for all 12 datasets.} We found that HiCOT's self-fit vocabulary path (used for every plain, non-\texttt{hicot\_*} dataset) had no stopword filter at all, unlike its own official \texttt{hicot\_*} vocabulary files and unlike every other baseline we checked (FASTopic, LDA already excluded stopwords unconditionally; BERTopic had the identical gap, also fixed). The corrected re-run changed HiCOT's own numbers only modestly, and does not change the paper's central classification finding.
\item \textbf{No genuine same-run comparison exists between VAE-BM and every baseline.} Every cluster and classification number in this paper is stitched from separate sweeps (different launch scripts, sometimes different compute environments) rather than one run containing every model; we report this explicitly rather than implying a single controlled comparison.
\item \textbf{A reproducible collapse in the HiCOT/FASTopic baselines on GoogleNews.} Both produce near-chance accuracy on the GoogleNews-family datasets in both the cluster and classification experiments, on two independent compute environments -- consistent with a shared degenerate-assignment failure mode rather than two genuine, independent weaknesses. We flag this rather than presenting it as a clean win, and have not root-caused it.
\item \textbf{Classification protocol.} Our classification sweeps use a random stratified 80/20 split rather than HiCOT's own official held-out files; no $K\in\{50,100\}$ classification result exists, since the random-split path does not take a requested $K$.
\item \textbf{Seed count differs by experiment, as a documented convention.} Cluster uses one seed throughout ($k$-means with a fixed initialization is deterministic here). Classification uses 5 seeds for VAE-BM (ours), since which documents land in train vs.\ test genuinely varies by seed, but a single seed (42) for every baseline, matching this task's own established single-seed baseline convention -- we did not have the budget to extend every baseline to 5 seeds, and report this asymmetry rather than fabricating uncertainty for a single-run baseline.
\item \textbf{Cross-hardware, not cross-validated.} Some baseline results exist on two different GPU environments (FutureLab, labuai); we report both rather than averaging, but neither is a repeated-seed estimate of the other.
\item \textbf{Appendix ablations use a different embedder and a smaller dataset set.} The Product-of-Experts and jointly-trained-clustering ablations in Appendix~\ref{sec:appendix-ablation} use \texttt{all-MiniLM-L6-v2} on HiCOT's five benchmark datasets only, not the locked \texttt{gte-large}/\texttt{bge-large} configuration or the 12-dataset suite -- they are exploratory evidence about the fusion and clustering \emph{mechanisms}, not a like-for-like comparison against this paper's main VAE-BM (ours) results.
\end{itemize}

\bibliography{anthology,custom,vaebmpoe}
\bibliographystyle{acl_natbib}

\appendix
\section{Results Traceability}
\label{sec:appendix}
Every numerical result in this paper is traceable to a specific source file via \texttt{paper\_data/\allowbreak results\_manifest.json}, produced by \texttt{scripts/\allowbreak build\_paper\_results.py} from the raw sweep checkpoints on both compute environments, and rendered into the tables above by \texttt{scripts/\allowbreak build\_paper\_tables.py}. Each manifest record carries: experiment, model, dataset, $K$ (if applicable), seed, every metric value, status, source file path, compute environment, sweep name, and checkpoint-selection flag.

\section{Exploratory Fusion and Clustering Ablations}
\label{sec:appendix-ablation}
This appendix reports two exploratory variants that do \emph{not} appear as our proposed model anywhere in the main text: an uncertainty-weighted encoder fusion, and a jointly-trained clustering objective. Both were evaluated on HiCOT's five benchmark datasets under an older configuration (\texttt{all-MiniLM-L6-v2}, not the locked \texttt{gte-large}/\texttt{bge-large} setup used for every main-text VAE-BM (ours) result) and neither displaced $\alpha=0$ fixed fusion (Equation~\ref{eq:final-fusion}) as this paper's proposed configuration.

\subsection{Product-of-Experts fusion}
Instead of the fixed scalar $\alpha$ of Equation~\ref{eq:alpha-fusion}, a Product-of-Experts (PoE) combination~\citep{hinton2002poe} fuses the two branch posteriors, together with the standard-normal prior, by precision:
\begin{equation}
q_{\mathrm{PoE}}(\mathbf{z}\mid\mathbf{x}_b,\mathbf{x}_e) \;\propto\; p(\mathbf{z})\,q_b(\mathbf{z}\mid\mathbf{x}_b)\,q_e(\mathbf{z}\mid\mathbf{x}_e).
\label{eq:poe-def}
\end{equation}
For diagonal Gaussians this product is itself Gaussian and closed-form. Writing each branch's precision as $\boldsymbol{\tau}_b=\boldsymbol{\sigma}_b^{-2}$ and $\boldsymbol{\tau}_e=\boldsymbol{\sigma}_e^{-2}$, the fused mean is the precision-weighted average of the two branch means and the prior's zero mean:
\begin{equation}
\boldsymbol{\mu}_{\mathrm{PoE}} = \frac{\boldsymbol{\tau}_b\,\boldsymbol{\mu}_b+\boldsymbol{\tau}_e\,\boldsymbol{\mu}_e}{\mathbf{1}+\boldsymbol{\tau}_b+\boldsymbol{\tau}_e}.
\label{eq:poe-mean}
\end{equation}
Whichever branch reports lower variance (higher confidence) for a given document and latent dimension dominates Equation~\ref{eq:poe-mean} for exactly that document and dimension, automatically and without a tuned hyperparameter -- in contrast to Equation~\ref{eq:alpha-fusion}'s $\alpha$, which is identical for every document in the corpus. We refer to this variant as VAE-BM-PoE.

\subsection{Jointly-trained clustering}
As an alternative to the two-stage $k$-means of \S\ref{sec:dec}, VAE-BM-DEC trains cluster centroids $\mathbf{c}_1,\dots,\mathbf{c}_K$ jointly with the encoder/decoder, following Deep Embedded Clustering~\citep{xie2016dec}. A Student-$t$ kernel gives each point's soft assignment to centroid $k$:
\begin{equation}
q_{ik} = \frac{\left(1+\|\mathbf{z}_i-\mathbf{c}_k\|^2\right)^{-1}}{\sum_j\left(1+\|\mathbf{z}_i-\mathbf{c}_j\|^2\right)^{-1}}.
\label{eq:dec-q}
\end{equation}
This is sharpened once per epoch into a target distribution that up-weights high-confidence assignments,
\begin{equation}
p_{ik} = \frac{q_{ik}^2/f_k}{\sum_j q_{ij}^2/f_j}, \qquad f_k=\sum_i q_{ik},
\label{eq:dec-p}
\end{equation}
and an additional loss term, weighted by $\lambda_c=0.1$, penalizes the KL divergence between target and soft assignment:
\begin{equation}
\mathrm{KL}(P\|Q) = \sum_i\sum_k p_{ik}\log\frac{p_{ik}}{q_{ik}}.
\label{eq:dec-kl}
\end{equation}
Centroids are initialized via $k$-means on the untrained encoder's own $\boldsymbol{\mu}$ before the first optimization step. Final cluster assignment is $y_d=\arg\max_k q_{dk}$, not a separate $k$-means call.

\subsection{Ablation result}
Table~\ref{tab:ablation} compares fixed-$\alpha$ fusion ($\alpha=0.99$, VAE-BM's own original default -- this ablation isolates the fusion \emph{mechanism} at the weight VAE-BM was originally proposed with, separately from the $\alpha$-value finding of \S\ref{sec:encoder}), VAE-BM-PoE, and VAE-BM-DEC on HiCOT's five benchmark datasets. VAE-BM-PoE and VAE-BM-DEC's cluster-experiment checkpoints were selected using the same ground-truth labels the experiment ultimately reports accuracy against (an oracle, label-informed selection, since the cluster experiment is transductive and has no held-out split to select on instead); the fixed-$\alpha$ row uses no such selection. Uncertainty-weighted fusion (PoE) improves ACC over fixed-$\alpha$ fusion on 4 of 5 datasets (all but IMDB-H); jointly-trained clustering (DEC) improves ACC on only 2 of 5 (20NG-H, IMDB-H), and both improvements are within 0.004 ACC of fixed-$\alpha$'s own value.

\begin{table}[t]
\centering
\small
\caption{Cluster ACC/NMI, HiCOT's five benchmark datasets, \texttt{all-MiniLM-L6-v2} embedder. VAE-BM-PoE and VAE-BM-DEC use oracle-informed checkpoint selection (this appendix, above); the fixed-$\alpha=0.99$ row does not. \textbf{Bold} = best per dataset+metric.}
\label{tab:ablation}
\adjustbox{max width=\columnwidth}{
\begin{tabular}{lcccccccccc}
\toprule
 & \multicolumn{2}{c}{20NG-H} & \multicolumn{2}{c}{AGNews-H} & \multicolumn{2}{c}{GNews-H} & \multicolumn{2}{c}{IMDB-H} & \multicolumn{2}{c}{SearchSnip.-H} \\
\cmidrule(lr){2-3}\cmidrule(lr){4-5}\cmidrule(lr){6-7}\cmidrule(lr){8-9}\cmidrule(lr){10-11}
 & ACC & NMI & ACC & NMI & ACC & NMI & ACC & NMI & ACC & NMI \\
\midrule
Fixed $\alpha{=}0.99$ & 0.549 & 0.571 & 0.815 & 0.541 & 0.674 & \textbf{0.866} & 0.790 & 0.258 & 0.658 & 0.419 \\
VAE-BM-PoE & \textbf{0.624} & \textbf{0.601} & \textbf{0.852} & \textbf{0.609} & \textbf{0.708} & 0.861 & 0.757 & 0.201 & \textbf{0.812} & \textbf{0.614} \\
VAE-BM-DEC & 0.553 & 0.545 & 0.737 & 0.450 & 0.509 & 0.717 & \textbf{0.793} & \textbf{0.265} & 0.537 & 0.323 \\
\bottomrule
\end{tabular}
}
\end{table}

\section{Full Clustering Metrics by Dataset}
\label{sec:appendix-cluster-full}
Tables~\ref{tab:cluster-full-google_news_t}--\ref{tab:cluster-full-pascal_flickr} report the remaining eight datasets of the 12-dataset suite (\S\ref{sec:cluster-results} shows the other four), in the identical full-11-metric, one-table-per-dataset format. Every baseline shown has a genuinely completed result on all 12 datasets.

Computed across all 12 datasets together with \S\ref{sec:cluster-results}'s four: \textbf{VAE-BM (ours) beats FASTopic on ACC and NMI on 12 of 12 datasets, beats HiCOT on ACC and NMI on 12 of 12 datasets, and beats BERTopic on ACC on 6 of 12 and NMI on 7 of 12} -- BERTopic, the one embedding-hybrid baseline, is again the closest competitor, consistent with \S\ref{sec:cluster-results}'s reading that VAE-BM (ours) trades most closely with baselines that themselves lean on embedding clustering.

\begin{table}[t]
\centering
\scriptsize
\caption{Full clustering metrics, GNews-T. Lower is better for DB only; higher is better for every other metric. \textbf{Bold}/\underline{underline} = best/second-best.}
\label{tab:cluster-full-google_news_t}
\adjustbox{max width=\columnwidth}{
\begin{tabular}{lrrrrrrrrrrr}
\toprule
Model & ACC & NMI & ARI & AMI & Hom. & Comp. & V-m. & Pur. & Sil. & DB & CH \\
\midrule
FASTopic & 0.233 & 0.485 & -- & -- & -- & -- & -- & 0.292 & -- & -- & -- \\
HiCOT & 0.040 & 0.007 & -- & -- & -- & -- & -- & 0.040 & -- & -- & -- \\
BERTopic & \textbf{0.759}$^{\mathrm{LU}}$ & \textbf{0.908}$^{\mathrm{LU}}$ & -- & -- & -- & -- & -- & \textbf{0.908}$^{\mathrm{LU}}$ & -- & -- & -- \\
SBERT-MiniLM & 0.667 & 0.859 & -- & -- & -- & -- & -- & 0.844 & -- & -- & -- \\
SBERT-GTE & \underline{0.714} & \underline{0.884} & -- & -- & -- & -- & -- & \underline{0.873} & -- & -- & -- \\
\midrule
VAE-BM & 0.664 & 0.857 & -- & -- & -- & -- & -- & 0.848 & -- & -- & -- \\
\bottomrule
\end{tabular}
}
\end{table}

\begin{table}[t]
\centering
\scriptsize
\caption{Full clustering metrics, BBCNews. Lower is better for DB only; higher is better for every other metric. \textbf{Bold}/\underline{underline} = best/second-best.}
\label{tab:cluster-full-bbc_news}
\adjustbox{max width=\columnwidth}{
\begin{tabular}{lrrrrrrrrrrr}
\toprule
Model & ACC & NMI & ARI & AMI & Hom. & Comp. & V-m. & Pur. & Sil. & DB & CH \\
\midrule
LDA & 0.776$^{\mathrm{LU}}$ & 0.606$^{\mathrm{LU}}$ & 0.566$^{\mathrm{LU}}$ & 0.605$^{\mathrm{LU}}$ & 0.597$^{\mathrm{LU}}$ & 0.616$^{\mathrm{LU}}$ & 0.606$^{\mathrm{LU}}$ & 0.776$^{\mathrm{LU}}$ & \underline{0.588}$^{\mathrm{LU}}$ & \underline{0.600}$^{\mathrm{LU}}$ & \textbf{2947.52}$^{\mathrm{LU}}$ \\
GloCOM & 0.852 & 0.697 & \underline{0.708} & \underline{0.696} & \underline{0.688} & \underline{0.707} & \underline{0.697} & 0.852 & \textbf{0.659} & 0.711 & \underline{2721.06} \\
ECRTM & 0.733 & 0.581 & 0.521 & 0.581 & 0.581 & 0.582 & 0.581 & 0.733 & 0.576 & \textbf{0.587} & 2678.17 \\
S2WTM & 0.540 & 0.409 & 0.338 & 0.407 & 0.378 & 0.446 & 0.409 & 0.551 & 0.382 & 0.914 & 1696.05 \\
FASTopic & 0.704 & 0.560 & -- & -- & -- & -- & -- & 0.731 & -- & -- & -- \\
HiCOT & 0.924 & 0.786 & -- & -- & -- & -- & -- & 0.924 & -- & -- & -- \\
BERTopic & 0.924$^{\mathrm{LU}}$ & 0.795$^{\mathrm{LU}}$ & -- & -- & -- & -- & -- & 0.924$^{\mathrm{LU}}$ & -- & -- & -- \\
SBERT-MiniLM & 0.943 & 0.831 & -- & -- & -- & -- & -- & 0.943 & -- & -- & -- \\
SBERT-GTE & \underline{0.961} & \underline{0.876} & -- & -- & -- & -- & -- & \underline{0.961} & -- & -- & -- \\
\midrule
VAE-BM & \textbf{0.961}$^{\mathrm{LU}}$ & \textbf{0.877}$^{\mathrm{LU}}$ & \textbf{0.908}$^{\mathrm{LU}}$ & \textbf{0.877}$^{\mathrm{LU}}$ & \textbf{0.877}$^{\mathrm{LU}}$ & \textbf{0.877}$^{\mathrm{LU}}$ & \textbf{0.877}$^{\mathrm{LU}}$ & \textbf{0.961}$^{\mathrm{LU}}$ & 0.069$^{\mathrm{LU}}$ & 3.527$^{\mathrm{LU}}$ & 91.89$^{\mathrm{LU}}$ \\
\bottomrule
\end{tabular}
}
\end{table}

\begin{table}[t]
\centering
\scriptsize
\caption{Full clustering metrics, Tweet. Lower is better for DB only; higher is better for every other metric. \textbf{Bold}/\underline{underline} = best/second-best.}
\label{tab:cluster-full-tweet}
\adjustbox{max width=\columnwidth}{
\begin{tabular}{lrrrrrrrrrrr}
\toprule
Model & ACC & NMI & ARI & AMI & Hom. & Comp. & V-m. & Pur. & Sil. & DB & CH \\
\midrule
FASTopic & 0.477 & 0.637 & -- & -- & -- & -- & -- & 0.559 & -- & -- & -- \\
HiCOT & 0.161 & 0.152 & -- & -- & -- & -- & -- & 0.168 & -- & -- & -- \\
BERTopic & 0.604$^{\mathrm{LU}}$ & \underline{0.876}$^{\mathrm{LU}}$ & -- & -- & -- & -- & -- & \textbf{0.924}$^{\mathrm{LU}}$ & -- & -- & -- \\
SBERT-MiniLM & \textbf{0.648} & 0.875 & -- & -- & -- & -- & -- & 0.912 & -- & -- & -- \\
SBERT-GTE & \underline{0.640} & \textbf{0.881} & -- & -- & -- & -- & -- & \underline{0.922} & -- & -- & -- \\
\midrule
VAE-BM & 0.618 & 0.847 & -- & -- & -- & -- & -- & 0.876 & -- & -- & -- \\
\bottomrule
\end{tabular}
}
\end{table}

\begin{table*}[t]
\centering
\small
\caption{Document clustering, full metric set, StackOv.. Higher is better for every metric except Davies-Bouldin (DB), where lower is better. $^{\mathrm{LU}}$ marks a labuai value (shown only when no FutureLab value exists for that model/dataset). \textbf{Bold}/\underline{underline} = best/second-best per metric.}
\label{tab:cluster-full-stack_overflow}
\adjustbox{max width=\textwidth}{
\begin{tabular}{lrrrrrrrrrrr}
\toprule
Model & ACC & NMI & ARI & AMI & Hom. & Comp. & V-m. & Pur. & Sil. & DB & CH \\
\midrule
LDA & 0.305$^{\mathrm{LU}}$ & 0.219$^{\mathrm{LU}}$ & 0.149$^{\mathrm{LU}}$ & 0.217$^{\mathrm{LU}}$ & 0.218$^{\mathrm{LU}}$ & 0.220$^{\mathrm{LU}}$ & 0.219$^{\mathrm{LU}}$ & 0.319$^{\mathrm{LU}}$ & 0.141$^{\mathrm{LU}}$ & 1.536$^{\mathrm{LU}}$ & 1041.82$^{\mathrm{LU}}$ \\
GloCOM & 0.718 & 0.680 & 0.543 & 0.679 & 0.666 & 0.694 & 0.680 & 0.718 & \textbf{0.729} & \textbf{0.815} & \textbf{5645.58} \\
ECRTM & 0.275 & 0.188 & 0.056 & 0.185 & 0.160 & 0.228 & 0.188 & 0.277 & \underline{0.149} & \underline{1.277} & 907.33 \\
S2WTM & 0.169 & 0.150 & 0.054 & 0.148 & 0.139 & 0.164 & 0.150 & 0.172 & 0.125 & 1.339 & \underline{1985.31} \\
FASTopic & 0.255 & 0.240 & 0.125 & 0.237 & 0.235 & 0.245 & 0.240 & 0.269 & 0.029 & 2.932 & 548.92 \\
HiCOT & 0.375 & 0.355 & 0.254 & 0.353 & 0.337 & 0.375 & 0.355 & 0.379 & 0.148 & 1.606 & 1947.85 \\
BERTopic & \textbf{0.761} & \underline{0.708} & \underline{0.645} & \underline{0.707} & \underline{0.708} & \underline{0.708} & \underline{0.708} & \textbf{0.761} & 0.075 & 4.214 & 292.34 \\
\midrule
\textbf{VAE-BM (ours)} & \underline{0.725}$^{\mathrm{LU}}$ & \textbf{0.791}$^{\mathrm{LU}}$ & \textbf{0.665}$^{\mathrm{LU}}$ & \textbf{0.790}$^{\mathrm{LU}}$ & \textbf{0.781}$^{\mathrm{LU}}$ & \textbf{0.801}$^{\mathrm{LU}}$ & \textbf{0.791}$^{\mathrm{LU}}$ & \underline{0.760}$^{\mathrm{LU}}$ & 0.078$^{\mathrm{LU}}$ & 3.579$^{\mathrm{LU}}$ & 358.76$^{\mathrm{LU}}$ \\
\bottomrule
\end{tabular}
}
\end{table*}

\begin{table}[t]
\centering
\scriptsize
\caption{Full clustering metrics, Biomed.. Lower is better for DB only; higher is better for every other metric. \textbf{Bold}/\underline{underline} = best/second-best.}
\label{tab:cluster-full-biomedical}
\adjustbox{max width=\columnwidth}{
\begin{tabular}{lrrrrrrrrrrr}
\toprule
Model & ACC & NMI & ARI & AMI & Hom. & Comp. & V-m. & Pur. & Sil. & DB & CH \\
\midrule
LDA & 0.175$^{\mathrm{LU}}$ & 0.102$^{\mathrm{LU}}$ & 0.055$^{\mathrm{LU}}$ & 0.099$^{\mathrm{LU}}$ & 0.102$^{\mathrm{LU}}$ & 0.102$^{\mathrm{LU}}$ & 0.102$^{\mathrm{LU}}$ & 0.191$^{\mathrm{LU}}$ & 0.133$^{\mathrm{LU}}$ & 1.543$^{\mathrm{LU}}$ & 954.57$^{\mathrm{LU}}$ \\
GloCOM & 0.467 & 0.408 & \underline{0.283} & \underline{0.406} & \underline{0.395} & \textbf{0.422} & \underline{0.408} & 0.468 & \textbf{0.606} & \textbf{0.889} & \textbf{3260.47} \\
ECRTM & 0.180 & 0.128 & 0.020 & 0.125 & 0.098 & 0.185 & 0.128 & 0.180 & 0.047 & \underline{1.054} & 564.44 \\
S2WTM & 0.170 & 0.141 & 0.051 & 0.138 & 0.130 & 0.153 & 0.141 & 0.173 & \underline{0.154} & 1.275 & \underline{1753.12} \\
FASTopic & 0.266 & 0.203 & -- & -- & -- & -- & -- & 0.269 & -- & -- & -- \\
HiCOT & 0.309 & 0.264 & -- & -- & -- & -- & -- & 0.316 & -- & -- & -- \\
BERTopic & 0.453$^{\mathrm{LU}}$ & 0.387$^{\mathrm{LU}}$ & -- & -- & -- & -- & -- & 0.467$^{\mathrm{LU}}$ & -- & -- & -- \\
SBERT-MiniLM & \underline{0.506} & \underline{0.442} & -- & -- & -- & -- & -- & \textbf{0.522} & -- & -- & -- \\
SBERT-GTE & \textbf{0.508} & \textbf{0.460} & -- & -- & -- & -- & -- & \underline{0.517} & -- & -- & -- \\
\midrule
VAE-BM & 0.441$^{\mathrm{LU}}$ & 0.417$^{\mathrm{LU}}$ & \textbf{0.287}$^{\mathrm{LU}}$ & \textbf{0.415}$^{\mathrm{LU}}$ & \textbf{0.416}$^{\mathrm{LU}}$ & \underline{0.418}$^{\mathrm{LU}}$ & \textbf{0.417}$^{\mathrm{LU}}$ & 0.473$^{\mathrm{LU}}$ & 0.036$^{\mathrm{LU}}$ & 4.243$^{\mathrm{LU}}$ & 225.25$^{\mathrm{LU}}$ \\
\bottomrule
\end{tabular}
}
\end{table}

\begin{table}[t]
\centering
\scriptsize
\caption{Full clustering metrics, Banking77. Lower is better for DB only; higher is better for every other metric. \textbf{Bold}/\underline{underline} = best/second-best.}
\label{tab:cluster-full-banking77}
\adjustbox{max width=\columnwidth}{
\begin{tabular}{lrrrrrrrrrrr}
\toprule
Model & ACC & NMI & ARI & AMI & Hom. & Comp. & V-m. & Pur. & Sil. & DB & CH \\
\midrule
FASTopic & 0.205 & 0.470 & -- & -- & -- & -- & -- & 0.225 & -- & -- & -- \\
HiCOT & 0.035 & 0.060 & -- & -- & -- & -- & -- & 0.035 & -- & -- & -- \\
BERTopic & 0.612$^{\mathrm{LU}}$ & \underline{0.802}$^{\mathrm{LU}}$ & -- & -- & -- & -- & -- & 0.641$^{\mathrm{LU}}$ & -- & -- & -- \\
SBERT-MiniLM & \underline{0.631} & 0.778 & -- & -- & -- & -- & -- & \underline{0.661} & -- & -- & -- \\
SBERT-GTE & \textbf{0.667} & \textbf{0.813} & -- & -- & -- & -- & -- & \textbf{0.705} & -- & -- & -- \\
\bottomrule
\end{tabular}
}
\end{table}

\begin{table*}[t]
\centering
\small
\caption{Document clustering, full metric set, M10. Higher is better for every metric except Davies-Bouldin (DB), where lower is better. $^{\mathrm{LU}}$ marks a labuai value (shown only when no FutureLab value exists for that model/dataset). \textbf{Bold}/\underline{underline} = best/second-best per metric.}
\label{tab:cluster-full-m10}
\adjustbox{max width=\textwidth}{
\begin{tabular}{lrrrrrrrrrrr}
\toprule
Model & ACC & NMI & ARI & AMI & Hom. & Comp. & V-m. & Pur. & Sil. & DB & CH \\
\midrule
LDA & 0.306$^{\mathrm{LU}}$ & 0.139$^{\mathrm{LU}}$ & 0.109$^{\mathrm{LU}}$ & 0.137$^{\mathrm{LU}}$ & 0.140$^{\mathrm{LU}}$ & 0.137$^{\mathrm{LU}}$ & 0.139$^{\mathrm{LU}}$ & 0.340$^{\mathrm{LU}}$ & 0.240$^{\mathrm{LU}}$ & 1.188$^{\mathrm{LU}}$ & 1395.30$^{\mathrm{LU}}$ \\
GloCOM & 0.576 & 0.465 & 0.410 & 0.464 & 0.471 & 0.460 & 0.465 & 0.639 & \textbf{0.580} & \textbf{0.577} & \textbf{5791.97} \\
ECRTM & 0.209 & 0.044 & 0.017 & 0.042 & 0.036 & 0.058 & 0.044 & 0.212 & 0.208 & \underline{1.021} & 881.00 \\
S2WTM & 0.316 & 0.196 & 0.112 & 0.195 & 0.179 & 0.218 & 0.196 & 0.342 & 0.217 & 1.128 & \underline{2091.16} \\
FASTopic & 0.300 & 0.181 & 0.134 & 0.179 & 0.181 & 0.182 & 0.181 & 0.318 & 0.075 & 2.199 & 753.81 \\
HiCOT & 0.261 & 0.115 & 0.089 & 0.113 & 0.116 & 0.114 & 0.115 & 0.280 & \underline{0.247} & 1.138 & 1442.38 \\
BERTopic & \underline{0.651} & \underline{0.502} & \underline{0.461} & \underline{0.501} & \underline{0.499} & \underline{0.505} & \underline{0.502} & \underline{0.676} & 0.037 & 4.551 & 131.18 \\
\midrule
\textbf{VAE-BM (ours)} & \textbf{0.690}$^{\mathrm{LU}}$ & \textbf{0.545}$^{\mathrm{LU}}$ & \textbf{0.485}$^{\mathrm{LU}}$ & \textbf{0.544}$^{\mathrm{LU}}$ & \textbf{0.547}$^{\mathrm{LU}}$ & \textbf{0.542}$^{\mathrm{LU}}$ & \textbf{0.545}$^{\mathrm{LU}}$ & \textbf{0.737}$^{\mathrm{LU}}$ & 0.044$^{\mathrm{LU}}$ & 4.083$^{\mathrm{LU}}$ & 157.23$^{\mathrm{LU}}$ \\
\bottomrule
\end{tabular}
}
\end{table*}

\begin{table*}[t]
\centering
\small
\caption{Document clustering, full metric set, Pascal-Fl.. Higher is better for every metric except Davies-Bouldin (DB), where lower is better. $^{\mathrm{LU}}$ marks a labuai value (shown only when no FutureLab value exists for that model/dataset). \textbf{Bold}/\underline{underline} = best/second-best per metric.}
\label{tab:cluster-full-pascal_flickr}
\adjustbox{max width=\textwidth}{
\begin{tabular}{lrrrrrrrrrrr}
\toprule
Model & ACC & NMI & ARI & AMI & Hom. & Comp. & V-m. & Pur. & Sil. & DB & CH \\
\midrule
LDA & 0.161$^{\mathrm{LU}}$ & 0.131$^{\mathrm{LU}}$ & 0.046$^{\mathrm{LU}}$ & 0.120$^{\mathrm{LU}}$ & 0.126$^{\mathrm{LU}}$ & 0.137$^{\mathrm{LU}}$ & 0.131$^{\mathrm{LU}}$ & 0.174$^{\mathrm{LU}}$ & \underline{0.271}$^{\mathrm{LU}}$ & \underline{1.098}$^{\mathrm{LU}}$ & 460.81$^{\mathrm{LU}}$ \\
GloCOM & 0.312 & 0.342 & 0.156 & 0.335 & 0.323 & 0.364 & 0.342 & 0.330 & \textbf{0.471} & 1.142 & \textbf{929.24} \\
ECRTM & 0.119 & 0.087 & 0.009 & 0.078 & 0.062 & 0.145 & 0.087 & 0.119 & 0.261 & \textbf{0.958} & 618.18 \\
S2WTM & 0.182 & 0.192 & 0.073 & 0.182 & 0.167 & 0.225 & 0.192 & 0.186 & 0.018 & 1.415 & \underline{669.85} \\
FASTopic & 0.186 & 0.198 & 0.072 & 0.186 & 0.183 & 0.215 & 0.198 & 0.196 & 0.008 & 2.861 & 286.73 \\
HiCOT & 0.180 & 0.167 & 0.066 & 0.156 & 0.154 & 0.182 & 0.167 & 0.195 & 0.081 & 1.528 & 469.83 \\
BERTopic & \textbf{0.382} & \textbf{0.416} & \textbf{0.235} & \textbf{0.408} & \textbf{0.414} & \textbf{0.417} & \textbf{0.416} & \textbf{0.403} & 0.026 & 4.215 & 65.77 \\
\midrule
\textbf{VAE-BM (ours)} & \underline{0.354}$^{\mathrm{LU}}$ & \underline{0.399}$^{\mathrm{LU}}$ & \underline{0.190}$^{\mathrm{LU}}$ & \underline{0.391}$^{\mathrm{LU}}$ & \underline{0.393}$^{\mathrm{LU}}$ & \underline{0.406}$^{\mathrm{LU}}$ & \underline{0.399}$^{\mathrm{LU}}$ & \underline{0.394}$^{\mathrm{LU}}$ & 0.037$^{\mathrm{LU}}$ & 3.662$^{\mathrm{LU}}$ & 82.78$^{\mathrm{LU}}$ \\
\bottomrule
\end{tabular}
}
\end{table*}


\end{document}
